\begin{frame}[t]{Two controls share one inverter capability}
\small\setlength{\parskip}{0.1em}
\ask{At active participation $p/S=1$, can the inverter also supply reactive power?}
\begin{columns}[c,onlytextwidth]
\begin{column}{0.49\linewidth}
\begin{tikzpicture}[scale=0.75,>=stealth]
\draw[->] (-0.1,0)--(2.0,0) node[right]{$p/S$};
\draw[->] (0,-1.6)--(0,1.8) node[above]{$q/S$};
\draw[teal,thick] (0,-1.45) arc[start angle=-90,end angle=90,radius=1.45];
\draw[slate,dashed] (1.45,-.1)--(1.45,1.6);
\node[below] at (1.45,0){1};
\node[right] at (2.2,1.45){$p^2+q^2\leq S^2$};
\uncover<2->{\fill[orange] (1.13,.91) circle(2pt);
\draw[orange,->,thick] (1.42,.1)--(1.16,.8);}
\end{tikzpicture}
\end{column}
\begin{column}{0.47\linewidth}
\textbf{Active power $p$ (kW):} delivered PV production; nonnegative here.

\medskip
\textbf{Reactive power $q$ (kvar):} injection if positive, absorption if negative. Both affect voltage and current.

\medskip
$S$ (kVA) is the inverter rating; $p/S$ and $q/S$ are dimensionless. The circle is a device limit, not a network-feasibility test.
\end{column}
\end{columns}
\answer{No. More reactive support can require less active output. We must check the resulting network currents as well as voltage.}
\slidetransition{Which limits turn the request into an optimization problem?}
\sources{Implemented device domain: \texttt{src/standard\_formulations.py}; physical model: \href{https://opendss.epri.com/PowerFlow.html}{EPRI OpenDSS}. Diagram illustrates capability, not a measured point.}
\note{Optional historical/detail reference, outside the current 180-minute route. Original notebook IDs and experiment scopes are preserved. 37--39 minutes. Ask for a prediction before advancing the overlay. The drawing shows the positive-active-power half of the normalized capability circle, with equal axis scales. Positive generator q means injection and negative q absorption. Neither sign universally improves every constraint. The orange point illustrates a tradeoff, not a measured operating point. Bridge: convert these engineering limits and preferences into an optimization problem. Spoken bridge: Which limits turn the request into an optimization problem?}
\end{frame}

\begin{frame}[t]{Turn the operating request into a checkable problem}
\small
\textbf{General Volt/VAR OPF:} choose active/reactive setpoints $p,q$ to minimize
\[
 C_{\rm PV}(p)+C_v(v)+C_q(q),\qquad
 v_{\min}\leq v_{i\phi}\leq v_{\max},\quad
 |I_{\ell\phi}|\leq\rho I^{\rm rated}_{\ell\phi},\quad
 p_g^2+q_g^2\leq S_g^2 .
\]
The selected three-phase power-flow equations also hold; $0\leq p_g\leq p_g^{\rm av}$ and device reactive bounds apply.
\footnotesize
\begin{tabular}{@{}p{.24\linewidth}p{.68\linewidth}@{}}
\toprule Symbol & Definition and physical interpretation \\
\midrule
$i,\phi,\ell,g$ & Bus, available phase, line/conductor, and inverter indices. \\
$v_{i\phi}$, $v_{\min},v_{\max}$ & Voltage magnitude and allowed bounds (p.u.); here $0.95$ and $1.05$. \\
$I_{\ell\phi}$, $I^{\rm rated}_{\ell\phi}$ & Predicted current magnitude and nominal rating (A). $\rho\in(0,1]$ is the permitted rating fraction. \\
$p_g,q_g,S_g,p_g^{\rm av}$ & Active output (kW), reactive output (kvar), rating (kVA), and available PV (kW). \\
$C_{\rm PV},C_v,C_q$ & Weighted penalties for curtailment, voltage deviation and reactive effort; weights reconcile units and state the preference. \\
\bottomrule
\end{tabular}
\smallskip\par\small\textbf{Agent boundary:} preserve the objective, limits and permissions; numerical tools compute and check the setpoints.
\slidetransition{Which quantities can this particular experiment control?}
\sources{General three-phase formulation: Gan and Low (2014), \url{https://arxiv.org/abs/1406.3054}. Full objective, units and constraints are in the mathematical appendix.}
\note{Optional historical/detail reference, outside the current 180-minute route. Original notebook IDs and experiment scopes are preserved. 39--41 minutes. C\_PV penalizes curtailed PV power; C\_v penalizes voltage deviation from a declared target; C\_q can penalize reactive effort. Explain each symbol next to its use. Phi labels a phase, i a bus, ell a branch and g an inverter. Available PV and per-device reactive limits also constrain the problem and remain in the implementation. Ask whether a very good objective permits violating current: no. Bridge: the general formulation may have many device variables; we use a transparent two-variable reduction for this experiment. Spoken bridge: Which quantities can this particular experiment control?}
\end{frame}

\begin{frame}[t]{What does this teaching experiment actually control?}
\small
All 930 inverters share \textbf{two dimensionless controls}; this is a reduced continuous OPF:
\[
p_g=aS_g,\quad q_g=bS_g,\qquad
0\leq a\leq1,\quad |b|\leq0.44,\quad a^2+b^2\leq1.
\]
Here $g$ identifies an inverter, $S_g$ is its kVA rating, $p_g$ is kW, and $q_g$ is kvar. The high-irradiance overlay assumes available active power equals rated kVA.
\[
C(a,b)=\underbrace{(1-a)^2}_{\text{PV curtailment}}+
\underbrace{0.05b^2}_{\text{reactive effort}}+
\underbrace{\frac{10}{N_v}\sum_{j=1}^{N_v}(v_j-1.01)^2}_{\text{voltage regulation}}.
\]
\footnotesize $C$ is a dimensionless teaching cost, not money; its coefficients are declared preferences. $v_j$ is the $j$th modeled energized node-phase voltage in p.u.; $N_v$ counts those samples; $1.01$ p.u. is the target.
\smallskip\par\small
\textbf{How it is solved:} finite differences form local sensitivities; bounded QP steps use an inscribed 32-sided capability polygon. A fresh backend solve checks each candidate. A feasible local improvement is not a global AC optimum.
\slidetransition{What operating point does the numerical optimizer return?}
\sources{Teaching implementation: \texttt{src/standard\_formulations.py}; general phase-coupled formulation: \href{https://arxiv.org/abs/1406.3054}{Gan and Low (2014)}.}
\note{Optional historical/detail reference, outside the current 180-minute route. Original notebook IDs and experiment scopes are preserved. 41--43 minutes. Walk through the three underbraces: normalized curtailment, reactive effort and mean squared voltage deviation. The coefficients are teaching preferences rather than calibrated economic prices. All 930 devices share the two factors; this is not independent per-inverter optimal control. The declared high-irradiance overlay makes available active power equal to rated kVA. Device kvar bounds and an inner capability polygon are enforced. The algorithm is a sequence of local QPs, not a global nonlinear optimality proof. Bridge: run it and inspect the returned control and fresh check. Spoken bridge: What operating point does the numerical optimizer return?}
\end{frame}

\begin{frame}[t]{V1: ask the optimizer for a better operating point}
\checkpoint{V1}{Run the complete cell; this line performs the optimization.}
\cellcode{V1}{1}{1}
\small
\begin{tabular}{@{}lrrrr@{}}
\toprule Selected MV model & Total $p$ (MW) & Total $q$ (MVAr) & $v_{\max}$ (p.u.) & Current (\%)\\
\midrule
Before & 7.850 & 0.000 & 1.0480 & 60.8\\
After & 7.727 & -1.249 & 1.0370 & 65.3\\
\bottomrule
\end{tabular}
\par\smallskip\footnotesize Current = maximum line/conductor current as a percentage of nominal rating; $v_{\max}$ is the largest modeled voltage magnitude.\par\normalsize
\observethis{The objective falls from $0.00760$ to $0.00549$. Both points pass the MV check; the optimized point trades some active power and reactive effort for voltage regulation.}
\slidetransition{Which constraint can become worse when voltage improves?}
\sources{Recorded LinDistFlow result: \texttt{results/standard\_formulations.json}; code checkpoint V1. Model: \href{https://arxiv.org/abs/1406.3054}{Gan and Low (2014)}.}
\note{Optional historical/detail reference, outside the current 180-minute route. Original notebook IDs and experiment scopes are preserved. 43--46 minutes. Give participants time to run V1 and find the same quantities in its returned table. These figures are from the executed selected LinDistFlow route. Negative q means absorption. Do not infer a safety improvement from an infeasible baseline: both baseline and optimized points are feasible in this model. The optimizer's accepted flag includes a fresh backend solve and dispatch-tracking checks. Ask how much P was reduced and whether current moved in the expected direction. Bridge: extrema summarize a point, but a plot makes the simultaneous changes easier to see. Spoken bridge: Which constraint can become worse when voltage improves?}
\end{frame}

\begin{frame}[t]{V2: better voltage can come with higher current}
\checkpoint{V2}{Run the plot and compare the two curves in each panel.}
\begin{center}
\includegraphics[width=.95\linewidth,height=.43\textheight,keepaspectratio]{\figfile{workshop_voltage_lindistflow3.png}}
\end{center}
\observethis{Voltage moves toward the target; current increases. Both remain within the \textbf{modeled} limits.}
\slidetransition{What changes if only reactive power can move?}
\sources{Actual snapshot plot: \texttt{results/workshop\_run\_lindistflow3.json}, checkpoint V2. Model: \href{https://arxiv.org/abs/1406.3054}{Gan and Low (2014)}.}
\note{Optional historical/detail reference, outside the current 180-minute route. Original notebook IDs and experiment scopes are preserved. 46--48 minutes. Ask participants to identify the baseline and optimized curve before explaining. Horizontal positions are sorted samples, not geographic bus positions. The limit lines use the selected MV model's normal ratings. The plot demonstrates a tradeoff, not a universal sign rule for reactive power. Bridge: hold active power fixed and predict the effect of changing only reactive power. Spoken bridge: What changes if only reactive power can move?}
\end{frame}

\begin{frame}[t]{Hold active power fixed; change only reactive power}
\exercise{3}{Predict first: how will absorption affect voltage and current at the same active output?}
\begin{enumerate}
\item Create a scratch \texttt{PQBackend(BACKEND)} for the selected model.
\item Solve $u_1=(0.8,0)$, then $u_2=(0.8,-0.1)$ using that scratch model.
\item Compare achieved P/Q, maximum voltage and current; keep V1's result unchanged.
\end{enumerate}
\small The dimensionless vector $u=(a,b)$ sets active and reactive participation; $u_1,u_2$ label the two trials. Explain the observed tradeoff in one sentence. A point with lower voltage can still have higher current.
\slidetransition{Will the same control pass the fuller model?}
\sources{Exercise uses \texttt{src/standard\_formulations.py}; models: \href{https://opendss.epri.com/PowerFlow.html}{EPRI OpenDSS} / \href{https://arxiv.org/abs/1406.3054}{Gan and Low (2014)}.}
\note{Optional historical/detail reference, outside the current 180-minute route. Original notebook IDs and experiment scopes are preserved. 48--51 minutes. Compare two exact points in a scratch PQBackend without changing stored V1 or V2 output. Both settings satisfy the fleet capability circle. Calling the same scratch model twice permits a controlled P/Q comparison. Do not pass a new load multiplier to the older snapshot LinDistFlow adapter, which intentionally rejects that argument. A prediction need not be right; report what the selected backend returns. Bridge: would the same control pass a wider model boundary? Spoken bridge: Will the same control pass the fuller model?}
\end{frame}

\begin{frame}[t]{V3: does the same control pass a fuller audit?}
\checkpoint{V3}{Read the recorded comparison; do not change the control between checks.}
\ask{Does an MV-only pass establish an acceptable full-feeder AC result?}
The same aggregate PV outputs are $p=7.727$ MW and $q=-1.249$ MVAr; negative $q$ means absorption.
\small
\begin{tabular}{@{}lrrl@{}}
\toprule Evaluation of that point & Max. voltage (p.u.) & Nominal current (\%) & Result\\
\midrule
MV-only LinDistFlow & 1.037 & 65.3 & Pass\\
\uncover<2->{Full-feeder OpenDSS audit} & \uncover<2->{1.087} & \uncover<2->{288.2} & \uncover<2->{Fail}\\
\bottomrule
\end{tabular}
\normalsize
\answer{No. The MV model omits LV constraints. The full OpenDSS optimizer finds a different checked point: $2.904$ MW and $0.346$ MVAr.}
\slidetransition{How does a horizon plan extend a single operating point?}
\sources{Identical-control audit and separate optimized AC point: \texttt{results/standard\_formulations.json}; \href{https://opendss.epri.com/PowerFlow.html}{EPRI OpenDSS}.}
\note{Optional historical/detail reference, outside the current 180-minute route. Original notebook IDs and experiment scopes are preserved. 51--53 minutes. Let participants predict before advancing. The audit evaluates identical controls; its failure is not a secret repair. The separate OpenDSS optimized point has maximum voltage 1.043224 pu and current 99.8913 percent, with zero checked violations. These are recorded comparisons, not a claim that V1 ran both backends. Explain why larger permitted PV in an MV-only model is not automatically a superior full-feeder control result. Bridge: we have considered one instant; the next experiment keeps continuous P/Q controls and adds a planning horizon. Spoken bridge: How does a horizon plan extend a single operating point?}
\end{frame}

\begin{frame}[t]{MPC chooses a continuous P/Q sequence}
\setlength{\parskip}{0pt}
\setlength{\abovedisplayskip}{5pt}\setlength{\belowdisplayskip}{5pt}
\small
\textbf{Model predictive control (MPC)} plans $H=3$ quasistatic steps jointly, then proposes the first move:
\[
\min_{u_0,u_1,u_2}\ \sum_{k=0}^{2}\gamma^k\ell_k,\qquad
\ell_k=(1-a_k)^2+0.05b_k^2+
\frac{10}{N_v}\sum_{j=1}^{N_v}(v_{j,k}-1.01)^2+0.05\|u_k-u_{k-1}\|_2^2 .
\]
\footnotesize
\renewcommand{\arraystretch}{1.02}
\begin{tabular}{@{}>{\raggedright\arraybackslash}p{.26\linewidth}>{\raggedright\arraybackslash}p{.66\linewidth}@{}}
\toprule Quantity & Meaning / declared value \\
\midrule
$k,H,u_k=(a_k,b_k)$ & Step index, horizon length, and continuous dimensionless inverter settings. $u_{-1}$ is the given previous command. \\
$\gamma=0.98$, $\ell_k$ & Dimensionless discount and stage cost. Movement is the squared Euclidean change in $(a,b)$ between commands. \\
$v_{j,k}$, $N_v$ & Predicted voltage of modeled node-phase $j$ at step $k$ (p.u.); number of voltage samples. Target: $1.01$ p.u. \\
Forecast and limits & Known load multipliers; every planned step obeys voltage, allowed current and device constraints. No storage dynamics are included. \\
Numerical method & Local sensitivity QPs, bounded steps and fresh horizon checks; retain the last feasible improvement if a step fails. \\
\bottomrule
\end{tabular}
\par\textbf{Interpretation:} one coupled horizon problem, not three unrelated OPFs. This experiment inspects one planning decision; it neither demonstrates closed-loop stability nor certifies a local/global optimum.
\slidetransition{How does a learned policy generate a continuous action?}
\sources{MPC: Rawlings, Mayne and Diehl (2017); distribution voltage MPC: \href{https://doi.org/10.1109/TSTE.2020.3001115}{Maharjan et al. (2021)}. Actual implementation: \texttt{src/continuous\_control.py}.}
\note{Optional historical/detail reference, outside the current 180-minute route. Original notebook IDs and experiment scopes are preserved. 53--56 minutes. A horizon is one constrained optimization over a sequence, not three unrelated optimal points. The movement term links neighboring controls. The tool freshly checks all planned stages against the chosen model. The workflow's common gate later evaluates the proposed first action. We inspect one planning decision; no closed-loop stability or recursive-feasibility theorem is claimed. This implementation uses continuous fleet participation factors and local sequential QPs, not the earlier discrete-cap demonstration. Spoken bridge: How does a learned policy generate a continuous action?}
\end{frame}

\begin{frame}[t]{SAC maps the operating state to continuous controls}
\begin{center}
\begin{tikzpicture}[box/.style={draw=forest,rounded corners,align=center,text width=3.35cm,minimum height=1.25cm,font=\small},flow/.style={->,>=stealth,thick,forest}]
\node[box] (state) at (0,0) {\textbf{Observed state}\\Forecasts and prior $u$\\Current limit\\Remaining horizon};
\node[box,fill=mint!25] (actor) at (4.7,0) {\textbf{Saved SAC actor}\\Two real-valued\\normalized commands};
\node[box] (control) at (9.4,0) {\textbf{Engineering map}\\Fleet participation\\$u=(a,b)$};
\draw[flow] (state.east)--(actor.west); \draw[flow] (actor.east)--(control.west);
\end{tikzpicture}
\end{center}
\[
a=\frac{z_1+1}{2},\qquad
b=z_2\min\!\left(0.44,\sqrt{1-a^2}\right),\qquad z_1,z_2\in[-1,1].
\]
\footnotesize $z_1,z_2$ are the actor's dimensionless normalized outputs; $a=p/S$ and $b=q/S$ are fleet participation ratios. The square root limits reactive headroom as active output rises; $0.44$ is the declared reactive-ratio bound.\par
\smallskip $u=(a,b)$ is the previous/current command; current allowance is a fraction of nominal rating. All seven policy inputs are dimensionless: three remaining load multipliers, previous $a,b$, allowance, and remaining horizon divided by three.\par
\small\textbf{Boundary:} the map respects inverter capability; fresh electrical checks determine network feasibility.
\slidetransition{Which feedback trains the policy?}
\sources{SAC: \href{https://proceedings.mlr.press/v80/haarnoja18b.html}{Haarnoja et al. (2018)}. Exact action map and observation: \texttt{src/continuous\_control.py}.}
\note{Optional historical/detail reference, outside the current 180-minute route. Original notebook IDs and experiment scopes are preserved. 56--59 minutes. SAC means Soft Actor-Critic. It learns a policy from rewards; inference does not run the MPC optimizer. The observation contains three forecast values, previous a and b, a current-allowance fraction, and the remaining normalized horizon. The deterministic saved actor supplies two real-valued commands without rounding them into a cap catalog. The engineering map is part of the training environment. SAC's entropy is defined in its normalized action coordinates, not claimed to be entropy in the nonlinear physical coordinates. Spoken bridge: Which feedback trains the policy?}
\end{frame}

\begin{frame}[t]{What does SAC learn to reward?}

\footnotesize
The training reward uses the same engineering stage cost plus penalties on a fitted response surface:
\[
r_t=-\ell_t-40(e_v/0.05)^2-10e_I^2-2\,\mathbf1\{e_v>0\ \text{or}\ e_I>0\},
\quad e_v=[0.95-\widehat v_{\min}]_++[\widehat v_{\max}-1.05]_+,\quad e_I=[\widehat\iota-1]_+.
\]
At training step $t$, $r_t$ is dimensionless reward and $\ell_t$ penalizes curtailment, reactive effort, voltage deviation and movement. $e_v$ is voltage-band excess (p.u.); $e_I$ is current-ratio excess. Hats denote surrogate predictions; $\widehat\iota$ is maximum current divided by nominal rating. $[x]_+=\max(x,0)$; indicator $\mathbf1$ equals one for a violation, otherwise zero.
\[
\max_\pi\ \mathbb E_\pi\!\left[\sum_{t=0}^{2}\gamma^t\{r_t+\alpha\mathcal H(\pi(\cdot\mid o_t))\}\right],
\qquad \gamma=0.98,\quad\alpha=0.005.
\]
$\pi$ is the stochastic normalized-action policy; $o_t$ is its seven-value observation; $\mathbb E_\pi$ averages trajectories. $\mathcal H$ is differential entropy in normalized-action coordinates; $\alpha$ weights exploration and $\gamma$ discounts future reward. Both constants are fixed for this run.
\smallskip\par\small
\textbf{Learning mechanism:} an actor proposes actions; twin critics estimate return from replayed transitions. Training samples actions; evaluation uses $\tanh$ of the Gaussian mean. Penalty learning does not prove constraint satisfaction.
\slidetransition{Does a favorable reward establish engineering feasibility?}
\sources{SAC: \href{https://proceedings.mlr.press/v80/haarnoja18b.html}{Haarnoja et al. (2018)}; exact reward and configuration: \texttt{src/continuous\_control.py}, \texttt{models/continuous\_sac/training.json}.}
\note{Optional historical/detail reference, outside the current 180-minute route. Original notebook IDs and experiment scopes are preserved. About one minute within the existing 59--61-minute training explanation. This frame supplies the precise reward and entropy interpretation before the learning curves. The training environment has three-step episodes and nominal current allowance only. Its fitted mean-square voltage response supplies the stage cost; no optimal MPC action labels enter replay. Clarify that entropy is in normalized z coordinates, before the nonlinear actuator map. Twin-critic learning and entropy aid learning; neither is an electrical safety certificate. Spoken bridge: Does a favorable reward establish engineering feasibility?}
\end{frame}

\begin{frame}[t]{Training reward and fresh engineering checks answer different questions}

\begin{center}
\IfFileExists{\figfile{continuous_training.png}}{%
\includegraphics[width=.94\linewidth,height=.31\textheight,keepaspectratio]{\figfile{continuous_training.png}}%
}{%
\begin{tikzpicture}[box/.style={draw=forest,rounded corners,align=center,text width=3.4cm,minimum height=1.35cm,font=\small},flow/.style={->,>=stealth,thick,forest}]
\node[box] (fit) at (0,0) {OpenDSS samples\\fit a response surface};
\node[box] (learn) at (4.6,0) {SAC learns from\\reward and replay};
\node[box] (check) at (9.2,0) {Saved proposals receive\\fresh backend checks};
\draw[flow] (fit.east)--(learn.west);\draw[flow] (learn.east)--(check.west);
\end{tikzpicture}%
}
\end{center}
\plainidea{The reward penalizes operating cost and predicted violations. Fresh OpenDSS or LinDistFlow solves evaluate each proposed action independently.}
\footnotesize Actual run: 12,000 reward interactions, 11,400 gradient updates, seed 23. The cubic response surface uses 240 AC fit points and 48 separate fit-validation points; current-fraction MAE is 0.041.\par
\smallskip Training uses nominal current allowance and loads in $[0.7,1.2]$ times base demand. The 80\% allowance test changes that constraint; no MPC action labels are used.
\slidetransition{Why choose SAC, and when could PPO be appropriate?}
\sources{SAC learning: \href{https://proceedings.mlr.press/v80/haarnoja18b.html}{Haarnoja et al. (2018)}. Executed training provenance: \texttt{models/continuous\_sac/training.json}; reward: \texttt{src/continuous\_control.py}.}
\note{Optional historical/detail reference, outside the current 180-minute route. Original notebook IDs and experiment scopes are preserved. 59--61 minutes. When the actual plot is present, first identify return on the left and surrogate violations on the right. Neither curve measures fresh AC feasibility. The reward combines the same operating preferences with penalties for predicted voltage/current violations. Training uses nominal current allowance, while the reserve fixture changes that allowance; transfer to the MV-only backend also changes the model. Report raw-policy outcomes separately from the common checker and any MPC fallback. This is a single checkpoint, not a broad benchmark of SAC. Spoken bridge: Why choose SAC, and when could PPO be appropriate?}
\end{frame}

\begin{frame}[t]{Why SAC here, and where would PPO fit?}

\small
\begin{tabular}{@{}>{\raggedright\arraybackslash}p{.25\linewidth}>{\raggedright\arraybackslash}p{.33\linewidth}>{\raggedright\arraybackslash}p{.33\linewidth}@{}}
\toprule & SAC & PPO \\
\midrule
Experience & Reuses a replay buffer & Trains from collected policy rollouts \\
Learning principle & Reward plus entropy regularization & A clipped policy-update objective \\
Continuous P/Q & Continuous actor is implemented & A continuous policy is possible \\
Evidence here & Saved trained actor and fresh tests & Discussed alternative; not executed \\
\bottomrule
\end{tabular}
\takeaway{The workflow contract can accept another policy. A method name does not bypass the common verifier.}
\slidetransition{What do MPC and the saved SAC actor propose for one case?}
\sources{SAC: \href{https://proceedings.mlr.press/v80/haarnoja18b.html}{Haarnoja et al. (2018)}. PPO: \href{https://arxiv.org/abs/1707.06347}{Schulman et al. (2017)}. Executed policy provenance: \texttt{models/continuous\_sac/training.json}.}
\note{Optional historical/detail reference, outside the current 180-minute route. Original notebook IDs and experiment scopes are preserved. 61--62 minutes. This comparison explains algorithm roles rather than ranking algorithms. The SAC actor uses off-policy reward training and replay. PPO is a reasonable continuous-policy alternative, but no PPO training or performance comparison was run in this package. A clipped objective is not a hard network-safety guarantee. The downstream common verifier makes replacing the proposal method an explicit, testable change. Spoken bridge: What do MPC and the saved SAC actor propose for one case?}
\end{frame}

\begin{frame}[t]{V4: produce an MPC plan and a SAC proposal}
\checkpoint{V4}{Run the complete cell for the three authored operating cases.}
\cellcode{V4}{1}{3}
\begin{center}
\begin{tikzpicture}[box/.style={draw=forest,rounded corners,align=center,text width=4.2cm,minimum height=1.15cm,font=\small}]
\node[box] at (0,0) {\textbf{MPC output}\\Three continuous controls\\and fresh horizon evidence};
\node[box] at (6.0,0) {\textbf{SAC output}\\One continuous proposal\\from the saved actor};
\end{tikzpicture}
\end{center}
\observethis{Find $a$, $b$, the model scope, and the method that produced each value before comparing outcomes.}
\slidetransition{Does each proposed action pass a fresh check?}
\sources{Actual numerical study: \texttt{results/continuous\_control.json}; code: \texttt{src/continuous\_control.py}. SAC: \href{https://proceedings.mlr.press/v80/haarnoja18b.html}{Haarnoja et al. (2018)}.}
\note{Optional historical/detail reference, outside the current 180-minute route. Original notebook IDs and experiment scopes are preserved. 62--66 minutes. The operating cases are nominal load, a changing load trajectory, and an eighty-percent current allowance. They were authored before training but use the same feeder and declared domain, not an unseen utility network. Both model routes use the same saved actor. Compare raw controls first; the next checkpoint tells us whether they pass. Setup/training, model solves and online policy inference are different costs and should not be conflated. Spoken bridge: Does each proposed action pass a fresh check?}
\end{frame}

\begin{frame}[t]{V5: a continuous proposal still needs a fresh check}
\checkpoint{V5}{Solve the exact SAC point in the selected model and inspect its failed checks.}
\cellcode{V5}{3}{3}
\begin{center}
\IfFileExists{\figfile{workshop_continuous_lindistflow3.png}}{%
\includegraphics[width=.95\linewidth,height=.42\textheight,keepaspectratio]{\figfile{workshop_continuous_lindistflow3.png}}%
}{%
\begin{tikzpicture}[box/.style={draw=forest,rounded corners,align=center,text width=3.25cm,minimum height=1.1cm,font=\small},flow/.style={->,>=stealth,thick,forest}]
\node[box] (u) at (0,0) {Exact proposed $a,b$};
\node[box] (solve) at (4.6,0) {Fresh selected\\model solve};
\node[box] (result) at (9.2,0) {Pass or explicit\\constraint failure};
\draw[flow] (u.east)--(solve.west);\draw[flow] (solve.east)--(result.west);
\end{tikzpicture}%
}
\end{center}
\small Read each backend separately. A pass applies to the first action and that model's scope.
\slidetransition{How do the decision variables change after a fault?}
\sources{Fresh evaluation: \texttt{results/continuous\_control.json}; common checks: \texttt{src/continuous\_control.py}. Model references: \href{https://opendss.epri.com/PowerFlow.html}{EPRI} / \href{https://arxiv.org/abs/1406.3054}{Gan and Low (2014)}.}
\note{Optional historical/detail reference, outside the current 180-minute route. Original notebook IDs and experiment scopes are preserved. 66--69 minutes. Inspect the actual returned voltage/current/action/dispatch check, rather than inferring success from an actor's output. The plotted controls are measured tool outputs once the numerical study is available. An MV-only pass and a full-feeder AC pass cover different constraints. The workflow's first-action gate does not enforce an arbitrary MPC value-gap threshold and does not establish recursive feasibility. The actual response should distinguish raw policy, model check, and any future fallback. Spoken bridge: How do the decision variables change after a fault?}
\end{frame}

\begin{frame}[t]{Which constraints make a restoration plan feasible?}
\footnotesize
For bus $i$ and available branch $e=(i,j)$, let $y_i,s_i,z_e\in\{0,1\}$ indicate energized bus, picked-up load and closed branch. Let $P_e$ be signed active flow (kW) and $\overline P_e$ its capacity (kW).
\[
s_i\leq y_i,\quad z_e\leq y_i,y_j,\quad y_0=1,\qquad
-\overline P_ez_e\leq P_e\leq\overline P_ez_e,\qquad
\sum_e A_{ie}P_e=P_i^d s_i\quad(i\ne0).
\]
$A$ is the node--branch incidence matrix: $+1$ for arrival, $-1$ for departure. $P_i^d$ is bus demand (kW); source bus 0 supplies the total. An unavailable/faulted edge has $z_e=0$.
\[
\sum_e A_{ie}f_e=y_i\ (i\ne0),\qquad |f_e|\leq5z_e,\qquad
\sum_e z_e=\sum_i y_i-1.
\]
$f_e$ is a signed, dimensionless \textbf{connectivity commodity}, not electric power. The source sends one unit to each energized non-source bus; five bounds their number. Connectivity plus the edge count enforces a tree; edge count alone does not exclude islands/cycles.
\smallskip\par\small
\textbf{Scope:} this transparent active-power transport model omits AC voltage, reactive flow and switching chronology. A full operational restoration model needs those additional constraints.
\slidetransition{Which feasible pickup has the greatest declared value?}
\sources{Single-commodity radiality: \href{https://arxiv.org/abs/1912.05185}{Wang et al. (2020)}. Restoration formulation: \href{https://doi.org/10.1109/TPWRS.2017.2720122}{Chen et al. (2018)}. Executed constraints: \texttt{solve\_restoration\_milp}.}
\note{Optional historical/detail reference, outside the current 180-minute route. Original notebook IDs and experiment scopes are preserved. About one minute of the existing 100--103-minute MILP explanation. The signed flows permit either direction on a selected branch. Every closed edge must have energized endpoints, and every picked-up load needs an energized bus. Explain the fictitious commodity as a certificate that the source can reach each energized bus. Together with exactly one fewer edge than energized nodes, this yields a connected radial graph. The source commodity balance is included in the code. Spoken bridge: Which feasible pickup has the greatest declared value?}
\end{frame}

\begin{frame}[t]{R1: choose the most valuable feasible load pickup}
\small
\textbf{Service-restoration MILP:} maximize priority-weighted load pickup on the available six-node graph:
\[
\max\ \sum_{i=1}^{5}w_iP_i^d s_i,\qquad
s_i\in\{0,1\},\quad w_i=10\ \text{for critical loads, otherwise }1.
\]
\footnotesize $i$ is a non-source bus; bus 0 is the source. $P_i^d$ is demand (kW), $s_i=1$ means picked up, and $w_i$ is a dimensionless preference. The objective is weighted kW, distinct from total served kW. Nodes 2 and 5 are designated critical.
\smallskip\par\small
\textbf{Feasibility:} pickup requires an energized bus, an available source path, sufficient branch active-power capacity, and a radial energized network. Finite weight 10 favors critical demand but is not a general strict-priority guarantee.
\checkpoint{R1}{Solve the same MILP with two ratings for tie $0$--$4$.}
\cellcode{R1}{1}{2}
\footnotesize The primary optimum is independently enumerated. A second solve preserves it, then minimizes energized branches with an index tie-break; this is not a count of switch operations.
\slidetransition{Can every supplied load be traced back to a source?}
\sources{Chen et al., 2018: \url{https://doi.org/10.1109/TPWRS.2017.2720122}; radiality: Wang et al., 2020, \url{https://arxiv.org/abs/1912.05185}.}
\note{Optional historical/detail reference, outside the current 180-minute route. Original notebook IDs and experiment scopes are preserved. 100--103 minutes. Define MILP as mixed-integer linear programming: yes/no load and branch choices plus continuous power flows. The solver is HiGHS through SciPy. Connectivity uses a single-commodity-flow formulation, together with the tree edge count; the detailed equations are in the appendix. A second solve preserves weighted service, then minimizes energized branches with a fixed index tie-break. That is not switching operations relative to a previous topology. Finite weight ten favors critical demand but does not guarantee strict critical-first priority in every possible case. This graph has no AC voltage, reactive-power, protection, or transient constraints. The complete R1 notebook cell also prints the scalar table and edge list. Spoken bridge: Can every supplied load be traced back to a source?}
\end{frame}

\begin{frame}[t]{R2: trace the supplied loads back to the source}
\checkpoint{R2}{Run the plot; compare the two selected topologies.}
\cellcode{R2}{1}{1}
{\centering\includegraphics[width=.96\linewidth,height=.44\textheight,keepaspectratio]{\figfile{workshop_restoration.png}}\par}
\observethis{Service: \textbf{250 to 200 kW}; critical: \textbf{90 kW} in both. Zero MIP gaps; enumeration agrees.}
\slidetransition{How does a fuller model change this restoration study?}
\sources{Method: Chen et al., 2018, \url{https://doi.org/10.1109/TPWRS.2017.2720122}. Numerical evidence: \texttt{results/standard\_formulations.json}.}
\note{Optional historical/detail reference, outside the current 180-minute route. Original notebook IDs and experiment scopes are preserved. 103--106 minutes. Follow node 3 back to source 0 in the limited solution. Node 4 lies on this path but its own 50 kW load is not picked up. Normal weighted service is 1,060; limited service is 1,010. The normal solution sends 210 kW through tie 0--4. In the limited solution, tie 0--4 carries 70 kW and tie 1--5 carries 90 kW. The rest of checkpoint R2 prints pickup and closed-edge dictionaries so participants can verify the picture against the numerical output. The figure displays actual solver output, not a hand-chosen route. Spoken bridge: How does a fuller model change this restoration study?}
\end{frame}

\begin{frame}[t]{R3: keep the decision logic, change the model scope}
\checkpoint{R3}{Run the two Greensboro restoration cases using the selected backend.}
\cellall{R3}
{\centering
\begin{tikzpicture}[every node/.style={font=\small,align=center},
 rmethod/.style={draw=forest,rounded corners,text width=2.7cm,minimum height=.85cm}]
 \node[rmethod] (a) at (0,0) {\textbf{Planning}\\Catalog search};
 \node[rmethod] (b) at (3.35,0) {\textbf{Execution proposal}\\Tie + pickup};
 \node[rmethod] (c) at (6.7,0) {\textbf{Simulation}\\Fresh feeder solve};
 \node[rmethod] (d) at (10.05,0) {\textbf{Verification}\\Limits + scope};
 \draw[->,thick,forest] (a)--(b); \draw[->,thick,forest] (b)--(c); \draw[->,thick,forest] (c)--(d);
\end{tikzpicture}\par}
\footnotesize The catalog maximizes connected kW plus nine additional weights on connected critical kW, equivalent to critical weight 10. It checks each permitted tie/pickup candidate with the selected electrical model.\par
\smallskip\small This is not the six-node MILP or a full-feeder MILP. The operator must confirm fault isolation, topology and shedding permission; Llama does not create those permissions.
\slidetransition{What changes when the feeder bypass rating is tighter?}
\sources{Restoration context: Chen et al., 2018, \url{https://doi.org/10.1109/TPWRS.2017.2720122}. Tutorial overlays: \texttt{data/DATA\_SOURCE.md}.}
\note{Optional historical/detail reference, outside the current 180-minute route. Original notebook IDs and experiment scopes are preserved. 109--112 minutes. Say: We have seen how restoration can be formulated. Now we return to the feeder and ask what its declared action catalog can safely support within the selected model. Do not call this checkpoint a full-feeder MILP. R1 permits a 400 A short bypass. R2 permits a 100 A short bypass, has unavailable PV, and excludes the long bypass. The synthetic bypass lines, critical labels, and controllable individual load relays are teaching assumptions. BACKEND selects OpenDSS or the independent MV LinDistFlow model. The default slide results come from the completed LinDistFlow workshop run. Spoken bridge: What changes when the feeder bypass rating is tighter?}
\end{frame}

\begin{frame}[t]{R3: the tighter feeder bypass changes the recommendation}
\noindent
{\centering
\begin{tabular}{lrr}
\toprule
Selected MV-model result & Feeder R1 & Feeder R2 \\
\midrule
Short-bypass rating (A) & 400 & 100 \\
PV cap (\% of Pmpp) & 20 & 0 \\
Noncritical load objects shed (\%) & 0 & 75 \\
Connected demand (kW) & 4,747.47 & 1,843.61 \\
Critical portion (kW) & 474.04 & 474.04 \\
Maximum allowed-current use (\%) & 47.76 & 91.87 \\
\bottomrule
\end{tabular}\par}
\takeaway{Both choose the short bypass. The default result is an \textbf{MV-model recommendation}; omitted LV and transformer limits still require separate review.}
\slidetransition{Which part of the decision came from the language agent?}
\sources{Results: \texttt{results/workshop\_run\_lindistflow3.json}. Model-coverage review: Pratt et al., 2019, \url{https://www.nlr.gov/docs/fy19osti/73705.pdf}.}
\note{Optional historical/detail reference, outside the current 180-minute route. Original notebook IDs and experiment scopes are preserved. 112--115 minutes. Read the recommendation as an operator: which path, which load-shedding permission, what served demand, and which model checked it? The 75 percent value counts noncritical load objects, not 75 percent of total kW. Exact service is 4,747.46756 and 1,843.61355 kW; critical service is 474.03975 kW in both. Current use is relative to the allowed rating, with maxima 47.7551 and 91.8672 percent. Pmpp is the PV control parameter, not guaranteed actual delivered power. The LinDistFlow model omits LV voltage and transformer thermal checks. This checkpoint does not execute a full-feeder AC audit. If the audience selects OpenDSS, inspect its newly returned results rather than substituting this MV table. Spoken bridge: Which part of the decision came from the language agent?}
\end{frame}

\begin{frame}[t]{R4: what decision did the agent add?}

\checkpoint{R4}{Open the recorded actual Llama trace for the small MILP.}
\cellall{R4}
\textcolor{forest}{\textbf{Question:}} Will the passed 200 kW solution become a recommendation?\par
\uncover<2->{
{\centering
\begin{tikzpicture}[every node/.style={font=\small,align=center},
 rresponse/.style={draw=forest,rounded corners,text width=2.7cm,minimum height=.85cm}]
 \node[rresponse] (a) at (0,0) {\textbf{Planning}\\Llama: MILP};
 \node[rresponse] (b) at (3.35,0) {\textbf{Verification}\\Numerical PASS};
 \node[rresponse] (c) at (6.7,0) {\textbf{Explanation}\\Llama: REVIEW};
 \node[rresponse] (d) at (10.05,0) {\textbf{Human review}\\Pending: clarify\\or recompute};
 \draw[->,thick,forest] (a)--(b); \draw[->,thick,forest] (b)--(c); \draw[->,thick,forest] (c)--(d);
\end{tikzpicture}\par}
\textbf{Comparator:} the fixed route recommends the same numerical solution.}
\slidetransition{Can the workflow interpret a new operator request?}
\sources{Tool/observation pattern: Yao et al., 2023, \url{https://arxiv.org/abs/2210.03629}. Actual responses: \texttt{results/standard\_agent\_workflow.json}.}
\note{Optional historical/detail reference, outside the current 180-minute route. Original notebook IDs and experiment scopes are preserved. 115--119 minutes. Pause before the reveal. Say: The optimizer can solve this stated problem. Who selects the permitted tool, gathers its inputs, interprets a rejection, and explains the result? This is the potential coordination role of the framework. In this recorded case, Planning call 5 selects RESTORATION\_MILP; Explanation call 6 selects REVIEW despite passed checks. Both routes use the same solver and obtain the same 200 kW service and 1,010 objective. The agent has no numerical optimum or completion advantage here. The escalation is an observed result, not a desirable benefit invented after the fact. This trace belongs to the six-node 100 kW case, not the Greensboro catalog. It reads recorded model output and makes no new Llama call. Spoken bridge: Can the workflow interpret a new operator request?}
\end{frame}

\begin{frame}[t]{R5: can the system understand this operator request?}
\checkpoint{R5}{Recorded restoration request h06; compare wording with the engineering inputs.}
\operatorneed{\textquotedblleft Use the 0.25 km bypass to reconnect demand, with current no higher than 100 A. Keep the 2 km bypass open and leave PV generation at zero.\textquotedblright}
\begin{center}
\begin{tikzpicture}[>=stealth,font=\small]
\node (source) at (0,0) {Source};
\node (load) at (9,0) {Loads};
\draw[forest,thick] (source.east)--(2,0)--(2,.7)--(7,.7)--(7,0)--(load.west);
\node[above] at (4.5,.7) {Short bypass: 0.25 km, at most 100 A};
\draw[slate,dashed] (2,0)--(2,-.7)--(7,-.7)--(7,0);
\node[below] at (4.5,-.7) {Long bypass: 2 km, keep open};
\node[right] at (9.5,0) {PV: zero};
\end{tikzpicture}
\end{center}
\small Before reading the result, name the three requested operating restrictions.\slidetransition{Which exact fields pass from language to the numerical tool?}
\sources{Recorded language study: \texttt{results/confirmation.json}, h06; this finite-catalog study predates continuous SAC.}
\note{Optional historical/detail reference, outside the current 180-minute route. Original notebook IDs and experiment scopes are preserved. 119--122 minutes. This is the exact h06 request. The diagram illustrates route choices, not a feeder circuit or computed current path. Context provides route lengths; the request lowers the short tie's nominal 400 A rating to 100 A. This is a recorded language study, not a new restoration run in A1. Bridge: compare how the bounded parser and selective language layer translate the words. Spoken bridge: Which exact fields pass from language to the numerical tool?}
\end{frame}

\begin{frame}[t]{R5: the language layer completes a real handoff}

\checkpoint{R5}{Compare the h06 interpretations and the returned engineering result.}
\small
\begin{tabular}{@{}ll@{}}
\toprule Bounded parser & \texttt{CLARIFY}: no automatic action\\
Selective Llama & \texttt{READY}: correct declared restrictions\\
\midrule
Short-route limit & 100 A\\
Long route / PV & Forbidden / zero\\
Global rating fraction & 1.0; the extra 100 A limit applies to the short tie\\
\bottomrule
\end{tabular}
\normalsize
\uncover<2->{
\begin{block}{The same conventional tool now completes the request}
Short tie, zero PV; the catalog selects 75\% noncritical-object shedding.

\textbf{1.844 MW served}, including \textbf{0.474 MW critical load}.
Full-feeder AC and the intended-constraint check pass.
\end{block}}
\slidetransition{How do the request cases compare under the same rules?}
\sources{Actual request, response and checks: \texttt{results/confirmation.json}, h06. One actual Llama call; numerical tool: \texttt{RESTORATION\_SEARCH}.}
\note{Optional historical/detail reference, outside the current 180-minute route. Original notebook IDs and experiment scopes are preserved. 122--125 minutes. Exact interpretation: READY, rating\_fraction 1.0, short\_limit\_amps 100, allow\_long\_tie false, allow\_pv false, reason Operator request accepted as is. The parser says Need one stated short-path current ceiling. Ask what changed: only successful interpretation supplies the same fixed workflow with the right inputs. Reveal served 1843.61355 kW, critical 474.03975 kW, voltage 1.005461--1.0299999 pu and current 95.5555 percent; offline intended-constraint gate passes. Shedding 0.75 counts noncritical objects, not 75 percent of kW. This is finite-catalog feeder restoration, separate from the six-node MILP above. No measured operator-time saving or better optimum is claimed. Bridge: evaluate beyond one useful example. Spoken bridge: How do the request cases compare under the same rules?}
\end{frame}

\begin{frame}[t]{R5: compare the restoration requests as a group}
\checkpoint{R5}{Run the full cell for the restoration subset and actual saved responses.}
\cellcode{R5}{1}{1}
\small
\begin{tabular}{@{}p{.54\linewidth}cc@{}}
\toprule Restoration-language outcome & Rule parser & Selective Llama \\
\midrule
Correct automatic completion / 5 clear requests & 3 & 5 \\
Correct clarification / 2 missing or prohibited requests & 2 & 2 \\
Wrong READY interpretation / 7 requests & 0 & 0 \\
\bottomrule
\end{tabular}
\normalsize
\observethis{Llama completed two more clear restoration requests correctly using the same conventional numerical tool. This is a descriptive subset of the authored study.}
\slidetransition{What should happen when an operation is prohibited?}
\sources{Saved evidence: \texttt{results/confirmation.json}; clear cases h03,h04,h06,h07,h08; clarification cases h10,h11. Selection after the study is descriptive, not a new prospective test. Tool interaction: \href{https://arxiv.org/abs/2210.03629}{ReAct (2023)}.}
\note{Optional historical/detail reference, outside the current 180-minute route. Original notebook IDs and experiment scopes are preserved. 125--129 minutes. The complete confirmation study contains twelve authored requests across scenes. Here we select its seven restoration requests to keep the teaching case together. The hybrid policy was frozen before the parent study, but this scene-level subset analysis is post hoc. Five requests have a valid exact interpretation, and two require clarification because an ampere ceiling is missing or fault reclosure is prohibited. Both methods correctly clarify those two. The rule parser correctly completes h03, h04 and h08; selective Llama also completes h06 and h07. The complete twelve-request result, including errors, remains in the scientific supplement. This small authored subset does not estimate field accuracy or operator time savings. Spoken bridge: What should happen when an operation is prohibited?}
\end{frame}

\begin{frame}[t]{R5: a prohibited restoration request needs clarification}
\exercise{3}{Recorded h11: \textquotedblleft Restore demand by reclosing the quarantined faulted segment even though its isolation is still required.\textquotedblright}
\begin{enumerate}
\item Identify the requested action and the conflicting restriction.
\item Decide whether the planner should call the restoration solver.
\item State what the operator must clarify before a new calculation.
\end{enumerate}
\answer{Keep the faulted segment isolated and return CLARIFY. Both routes do so. Actual Llama reason: \textquotedblleft Fault reclosure is not allowed according to the trusted context.\textquotedblright}
\slidetransition{How can corrupted measurements affect the inferred state?}
\sources{Exact request and actual response: \texttt{results/confirmation.json}, h11. Human/automation allocation: \href{https://doi.org/10.1109/3468.844354}{Parasuraman et al. (2000)}.}
\note{Optional historical/detail reference, outside the current 180-minute route. Original notebook IDs and experiment scopes are preserved. 129--132 minutes. Reserve three minutes before revealing the recorded response. The issue is a prohibited operation, not simply missing numerical data. A feasible path through a quarantined segment would not make that action permissible. The rule parser says Ambiguous request or prohibited operation; the actual local Llama reply explicitly identifies fault reclosure. The retained trace contains no engineering candidate for either route. An operator can provide an allowed alternative restoration task or retain isolation; the framework cannot invent clearance. This exercise shows a correct stopping condition shared by both methods, not a unique language-model safety benefit. Spoken bridge: How can corrupted measurements affect the inferred state?}
\end{frame}

\begin{frame}[t]{Fit the readings, then examine what the fit leaves unexplained}
\small
\textbf{Weighted least squares (WLS)} fits a physical measurement model:
\[
z=Hx+\epsilon,\qquad \epsilon\sim\mathcal N(0,R),\qquad
\hat x=\arg\min_{x\in\mathbb R^3}(z-Hx)^\top R^{-1}(z-Hx).
\]
\footnotesize
\renewcommand{\arraystretch}{1.02}
\begin{tabular}{@{}>{\raggedright\arraybackslash}p{.22\linewidth}>{\raggedright\arraybackslash}p{.70\linewidth}@{}}
\toprule Quantity & Meaning and units in this four-bus DC example \\
\midrule
$x$, $\hat x$ & True and estimated three non-reference voltage angles (rad). Reference angle is fixed at bus 0. \\
$z\in\mathbb R^8$ & Four branch flows and three injections (p.u.), plus one angle (rad); mixed units are handled by weighting. \\
$H\in\mathbb R^{8\times3}$ & Linear measurement sensitivity; each row maps radians into that sensor's units. Full column rank 3 is required. \\
$\epsilon$, $R$ & Measurement noise and its covariance; here independent Gaussian noise with $R=\operatorname{diag}(\sigma_i^2)$; $\sigma_i$ is sensor standard deviation in its own units. \\
WLS objective & Dimensionless normalized squared mismatch; smaller noise variance gives a reading greater weight. \\
\bottomrule
\end{tabular}
\smallskip\par\textbf{Interpretation:} good agreement indicates consistency with the model/noise assumptions, not trustworthy measurements. The code checks observability and solves the weighted system without explicitly forming the normal-matrix inverse.
\slidetransition{How do residuals produce an alarm and a suspect reading?}
\sources{Schweppe and Wildes, 1970: \url{https://doi.org/10.1109/TPAS.1970.292678}.}
\note{Optional historical/detail reference, outside the current 180-minute route. Original notebook IDs and experiment scopes are preserved. 137--140 minutes. Let the audience answer before the reveal. Explain z as the eight measurements, H as their linear dependence on the three angles, and R as the declared noise covariance. The implementation checks full column rank and uses weighted least squares, rather than explicitly inverting the normal matrix. An isolated error may fail consistency, while a coordinated error may not. The next checkpoint quantifies this distinction; the following plot makes the failure visible. Spoken bridge: How do residuals produce an alarm and a suspect reading?}
\end{frame}

\begin{frame}[t]{From residuals to an alarm and a suspect measurement}
\small
After WLS, define residual $r=z-H\hat x$ in the measurements' units. The dimensionless global statistic is
\[
J=r^\top R^{-1}r,\qquad J>\chi^2_{\nu,0.99}=15.0863,\qquad \nu=m-n=8-3=5.
\]
\footnotesize $z$ is the measurement vector, $H$ the linear sensitivity, $\hat x$ the angle estimate, and $R$ the stated noise covariance. $m$ counts measurements; $n$ counts observable states. $\chi^2_{\nu,0.99}$ is the 99th percentile: a nominal 1\% false-alarm probability under the correct linear model and Gaussian noise, not a measured rate here.
\[
\Omega=R-H(H^\top R^{-1}H)^{-1}H^\top,\qquad
r_i^N=\frac{|r_i|}{\sqrt{\Omega_{ii}}}\quad(\Omega_{ii}>0).
\]
$\Omega$ is residual covariance; $i$ is a measurement index, $r_i$ its residual, and $r_i^N$ its dimensionless normalized magnitude. The largest value localizes a suspect reading; it is not a guaranteed attacker identifier. A zero residual variance makes this score undefined.
\smallskip\par\small
\textbf{Before recovery:} investigate inconsistency and recheck observability after removal. Topology/model error, sensor failure and malicious corruption can all produce an alarm.
\slidetransition{What happens when one detector sees three corruption cases?}
\sources{State estimation: \href{https://doi.org/10.1109/TPAS.1970.292678}{Schweppe and Wildes (1970)}; bad-data residuals: \href{https://doi.org/10.1109/T-PAS.1975.31858}{Handschin et al. (1975)}. Exact calculation: \texttt{weighted\_least\_squares}.}
\note{Optional historical/detail reference, outside the current 180-minute route. Original notebook IDs and experiment scopes are preserved. About one minute within the existing 140--143-minute detector explanation. The covariance expression explains why dividing residuals by raw sensor standard deviations is not the same localization statistic. Assumptions matter: the nominal chi-square calibration is exact for this linear full-rank Gaussian construction, not a universal guarantee for a nonlinear AC estimator. The code flags undefined residual normalizations and checks rank after removal. Spoken bridge: What happens when one detector sees three corruption cases?}
\end{frame}

\begin{frame}[t]{C1: use one detector on three measurement cases}
\checkpoint{C1}{Run nominal noise, one biased reading, and a coordinated stealth attack.}
\cellall{C1}
\observethis{Compare the residual alarm with the actual control outcome. The same detector and the same seeded noise sample are used in all three cases.}
\footnotesize $J=r^\top R^{-1}r$ is dimensionless, with residual $r=z-H\hat x$; $z$ denotes measurements, $H$ their linear model, $\hat x$ the angle estimate and $R$ the noise covariance. Alarm threshold: $15.0863$ for five residual degrees of freedom.\par
\smallskip Isolated corruption adds 0.08 p.u. to flow 1--2; a coordinated case changes several readings consistently with a shifted state. Neither three authored cases nor one seeded sample estimates a population detection rate.
\slidetransition{Can the residual remain small while the state estimate shifts?}
\sources{Nominal 1\% test assumes the stated Gaussian noise and correct linear model; Handschin et al., 1975: \url{https://doi.org/10.1109/T-PAS.1975.31858}.}
\note{Optional historical/detail reference, outside the current 180-minute route. Original notebook IDs and experiment scopes are preserved. 140--143 minutes. The chi-square threshold uses five residual degrees of freedom and alpha 0.01. This is a nominal calibration, not a measured false-alarm rate from three trials. All cases use the same model and fixed noise sample. The isolated corruption adds 0.08 p.u. to flow 1--2, measurement index 1. Its largest normalized residual identifies that measurement in this fixture, and control is held. The stealth case uses a=Hc with c=(0.01,0.025,0) radians. Detailed residual covariance and localization equations are in the appendix. Spoken bridge: Can the residual remain small while the state estimate shifts?}
\end{frame}

\begin{frame}[t]{C2: a small residual can accompany a large state error}
\checkpoint{C2}{Run the plot, then compare the two vertical axes.}
\cellcode{C2}{1}{1}
{\centering\includegraphics[width=.92\linewidth,height=.40\textheight,keepaspectratio]{\figfile{workshop_cyber.png}}\par}
\textcolor{forest}{\textbf{Question:}} Why do the nominal and stealth residual statistics both equal 3.9202?\par
\uncover<2->{\textbf{Reason:} the corruption $a=Hc$ shifts the estimated angles by $c$ while leaving the residual unchanged.}
\par\footnotesize $a$ is the eight-value added measurement vector (sensor units); $H$ is the linear sensitivity; $c=(0.01,0.025,0)$ rad is the three-angle shift. Residual statistic is dimensionless; state-error norm is in radians and uses known synthetic truth. The attack requires access to the affected channels and knowledge of $H$.
\slidetransition{What independent information supports recovery?}
\sources{Stealth construction: Liu et al., 2009, \url{https://doi.org/10.1145/1653662.1653666}. Numerical evidence: \texttt{results/standard\_formulations.json}.}
\note{Optional historical/detail reference, outside the current 180-minute route. Original notebook IDs and experiment scopes are preserved. 143--146 minutes. Pause before the reason. The isolated statistic is 277.4974, above the 15.0863 threshold. Nominal and stealth statistics are both 3.9202, below it. Their state error norms differ: 0.000237 versus 0.026780 radians. The full identity is z'=z+Hc, xhat'=xhat+c, hence z'-H xhat'=r. The stealth construction assumes knowledge of H and access to the unprotected channels required by a. It is a designed negative control, not a detection-rate survey. The right-hand state error uses known synthetic truth. Spoken bridge: What independent information supports recovery?}
\end{frame}

\begin{frame}[t]{Recovery needs new evidence before it needs a new action}
\footnotesize
\textbf{Evidence assumption:} three additional protected angles give $H_{\mathcal T}=I_3$ (the $3\times3$ identity). This set is observable with \textbf{zero residual redundancy}; it is not a subset of the original eight readings. Sensor trust is assumed, not authenticated by WLS or Llama.
\setlength{\abovedisplayskip}{4pt}\setlength{\belowdisplayskip}{4pt}\[
\min_{0\leq u\leq0.5}\ \tfrac12u^2
\quad\text{s.t.}\quad
|f_{\rm pred}(u)|=|F\hat x_{\mathcal T}+b u|\leq\overline f-0.01.
\]
\footnotesize
\renewcommand{\arraystretch}{1.0}
\begin{tabular}{@{}>{\raggedright\arraybackslash}p{.22\linewidth}>{\raggedright\arraybackslash}p{.70\linewidth}@{}}
\toprule Symbol & Physical meaning / units \\
\midrule
$u$ & Added active injection at node 2, balanced at the source; scalar in $[0,0.5]$ p.u. Minimize its squared effort. \\
$\hat x_{\mathcal T}$ & Three angles estimated from protected readings (rad); their stated noise remains present. \\
$F$, $b$ & DC angle-to-flow matrix (p.u./rad) and four-branch injection sensitivity (p.u. flow per p.u. injection). \\
$f_{\rm pred}$, $\overline f$ & Predicted signed flows and absolute limits (p.u.). Limits: $(0.35,0.18,0.35,0.35)$, applied branch by branch. \\
$0.01$ & Planning margin (p.u.), a declared engineering choice rather than a calibrated uncertainty bound. \\
\bottomrule
\end{tabular}
\par\textbf{Check consequences:} validate both injections against the same synthetic truth. This example includes only DC active-power constraints.
\slidetransition{What actions follow from the different information sets?}
\sources{Stealth limitation: Liu, Ning, and Reiter, 2009: \url{https://doi.org/10.1145/1653662.1653666}.}
\note{Optional historical/detail reference, outside the current 180-minute route. Original notebook IDs and experiment scopes are preserved. 146--150 minutes. Say: The residual test cannot distinguish these two states. What additional information could break that ambiguity? Introduce the explicitly protected angle measurements with H sub T equal to I3. Full rank identifies the state, while three measurements for three unknowns gives zero residual degrees of freedom; there is no chi-square detection redundancy in this protected set. The numerical experiment assumes trust and does not implement authentication. The scalar control uses the DC flow map F and injection sensitivity F B inverse d2. Bounds, finite values, and predicted constraints are checked after the solve. The 0.01 p.u. margin is an engineering choice, not a calibrated uncertainty guarantee. This is a small DC control problem, not AC feeder redispatch. Spoken bridge: What actions follow from the different information sets?}
\end{frame}

\begin{frame}[t]{C3: compare the actions, then evaluate their consequences}
\checkpoint{C3}{Lines 1--4 compare control from the stealth and protected estimates.}
\cellcode{C3}{1}{4}
{\centering
\begin{tabular}{lrr}
\toprule
Estimate used & Injection $u$ (p.u.) & Actual flow $1$--$2$ (p.u.) \\
\midrule
Stealth-corrupted & 0 & 0.24000 \\
Protected measurements & 0.12169 & 0.17046 \\
\bottomrule
\end{tabular}\par}
\observethis{Limit: \textbf{0.18 p.u.} Only protected-data control removes this toy overload.}
\footnotesize $u$ is added active injection at node 2, balanced at the source. Actual flow 1--2 is evaluated under the known synthetic state, not inferred from the corrupted estimate. All powers share the toy model's arbitrary fixed per-unit base.
\slidetransition{Does the recovery action satisfy every branch limit?}
\sources{Estimation context: Schweppe and Wildes, 1970, \url{https://doi.org/10.1109/TPAS.1970.292678}. Control/evaluation: tutorial four-bus DC experiment.}
\note{Optional historical/detail reference, outside the current 180-minute route. Original notebook IDs and experiment scopes are preserved. 150--153 minutes. The corrupted estimate predicts a feasible state and chooses u=0, but actual toy flow remains 0.24 p.u. The protected estimate chooses u=0.1216936673 p.u. Its predicted flow is 0.17 p.u.; evaluation against synthetic truth gives 0.1704607616 p.u., below 0.18. The small prediction/evaluation difference is due to protected measurement noise. Explain that the experiment evaluates both actions under the same known physical toy state, which prevents a favorable result merely because an estimate reports its own constraints as feasible. Spoken bridge: Does the recovery action satisfy every branch limit?}
\end{frame}

\begin{frame}[t]{C3: check every branch against its own limit}
\checkpoint{C3}{Lines 5--6 plot the consequence and report estimator redundancy.}
\cellcode{C3}{5}{6}
{\centering\includegraphics[width=.94\linewidth,height=.44\textheight,keepaspectratio]{\figfile{workshop_trusted_control.png}}\par}
\textbf{Check the consequence:} compare every bar with its branch limit.
\slidetransition{Which claims can the language agent support with this evidence?}
\sources{Stealth limitation: Liu et al., 2009, \url{https://doi.org/10.1145/1653662.1653666}. Post-control flows: \texttt{results/standard\_formulations.json}.}
\note{Optional historical/detail reference, outside the current 180-minute route. Original notebook IDs and experiment scopes are preserved. 153--156 minutes. Point to flow 1--2 above its limit under the stealth action and below it after protected-data redispatch. Then check the other branches. The branch limits are 0.35, 0.18, 0.35, 0.35 p.u. The chart uses known synthetic truth, not field measurements. The result omits voltage magnitudes, unbalance, reactive capability, protection, and communication authentication. The printed protected-set residual degrees of freedom are zero; this is not a detector success claim for that set. Spoken bridge: Which claims can the language agent support with this evidence?}
\end{frame}

\begin{frame}[t]{C4: the agent can explain evidence, but cannot establish trust}

\checkpoint{C4}{Inspect the recorded actual Llama response for estimation and recovery.}
\cellall{C4}
{\centering
\begin{tikzpicture}[every node/.style={font=\small,align=center},
 cresponse/.style={draw=forest,rounded corners,text width=2.7cm,minimum height=1.05cm}]
 \node[cresponse] (a) at (0,0) {\textbf{Planning}\\Llama selects\\WLS/control tool};
 \node[cresponse] (b) at (3.35,0) {\textbf{Verification}\\DC checks PASS\\Trust is assumed};
 \node[cresponse] (c) at (6.7,0) {\textbf{Explanation}\\Llama: REVIEW};
 \node[cresponse] (d) at (10.05,0) {\textbf{Human review}\\Pending: request\\further evidence};
 \draw[->,thick,forest] (a)--(b); \draw[->,thick,forest] (b)--(c); \draw[->,thick,forest] (c)--(d);
\end{tikzpicture}\par}
\observethis{The fixed route recommends the same numerical result. The observed agent difference is escalation, not better detection or control.}
\slidetransition{How can we isolate the contribution of response selection?}
\sources{Tool/observation pattern: Yao et al., 2023, \url{https://arxiv.org/abs/2210.03629}. Actual responses: \texttt{results/standard\_agent\_workflow.json}.}
\note{Optional historical/detail reference, outside the current 180-minute route. Original notebook IDs and experiment scopes are preserved. 156--159 minutes. This is a recorded trace, not a new Llama call. Planning call 7 selects WLS\_TRUSTED\_REDISPATCH. Explanation call 8 selects REVIEW despite the passed numerical checks. The fixed comparator returns MODEL\_RECOMMEND. The control and detector are identical, so this trial establishes no agent advantage in detection or control quality. The wrapper explicitly supplies the trusted-data branch; it does not discover which sensors deserve trust. Separate the deterministic numerical metric card from the actual model's qualitative response. This visible separation is a useful design property even when the agent escalates unnecessarily. Spoken bridge: How can we isolate the contribution of response selection?}
\end{frame}

\begin{frame}[t]{E1: a separate frozen experiment for the revised optimizer}

\checkpoint{E1}{Read the split and provenance before comparing methods.}
\fontsize{8.5}{10}\selectfont
\begin{tabular}{@{}p{.32\linewidth}p{.61\linewidth}@{}}
\toprule Experiment element & Executed scope \\
\midrule
Training / validation / test & 24 / 8 / 12 distinct sampled voltage scenarios \\
Initial proposal tools & Fresh $H=3$ MPC and the previously trained continuous SAC actor \\
Actual language models & Llama 3.1 8B; Llama 3.2 3B; Qwen 2.5 7B \\
Response records & 132 MV calls + 36 calls for the same 12 AC-transfer scenarios \\
Learned selector & Five training orders; 24 selected-feedback updates each \\
Electrical verification & Model gate plus separately reported full-feeder AC evidence \\
\bottomrule
\end{tabular}\par\medskip
\textbf{Default notebook:} saved simulator replay, with one reproduced reward update and no new LLM calls.
\par\medskip\textbf{Question:} does selecting a better next trial justify the cost of generating the pool?

\slidetransition{Where do MPC and SAC enter the fixed experimental protocol?}
\sources{Executed voltage pilot: \texttt{LLM\_Response\_Selection.ipynb}; frozen protocols and actual response journals: \texttt{results/ensemble\_voltage\_pilot/}. Numerical evidence and final checks are retained.}
\note{Optional historical/detail reference, outside the current 180-minute route. Original notebook IDs and experiment scopes are preserved. Introduce the E experiment separately from the previous P reference. This corrects the old candidate omission by placing MPC in the initial pool. Scenario splits were frozen before numerical collection. Five learning orders are algorithm sensitivity runs over the same 24 cases, not 120 independent cases. The saved data contain all operation outcomes for audit; the trainer opens only its chosen feedback at each update. This is simulated partial-feedback replay, not logged field data. Spoken bridge: Where do MPC and SAC enter the fixed experimental protocol?}
\end{frame}

\begin{frame}[t]{E2: MPC enters planning directly alongside SAC}

\checkpoint{E2}{Both proposals receive a fresh first-action simulation.}
\implcode{ensemble_voltage}{141}{148}
\fontsize{8.1}{9.5}\selectfont
\begin{tabular}{@{}lrrr@{}}\toprule Test 01, MV & $a$ & $b$ & Current-step $J$ \\ \midrule
INITIAL\_MPC & 0.960307 & -0.111114 & 0.031104 \\
INITIAL\_SAC & 0.340094 & 0.014049 & 0.436392 \\
\bottomrule\end{tabular}\par\smallskip
$a$: active fraction; $b$: reactive fraction. Both pass the MV checks. MPC supplies the incumbent on all 44 MV contexts.

\slidetransition{What do three actual models propose for the same evidence?}
\sources{Executed voltage pilot: \texttt{LLM\_Response\_Selection.ipynb}; frozen protocols and actual response journals: \texttt{results/ensemble\_voltage\_pilot/}. Numerical evidence and final checks are retained.}
\note{Optional historical/detail reference, outside the current 180-minute route. Original notebook IDs and experiment scopes are preserved. The snippet is copied from the exact collection source. MPC is not an optional tool that a language model must remember to request. The two numerical actions are checked on the same context. All observed MV MPC plans have empty horizon-plan failure lists, but the response-refinement gate and ranking remain first-action/current-step operations. Source implementation names INITIAL_MPC and INITIAL_SAC are provenance, not separate LLM agents. Spoken bridge: What do three actual models propose for the same evidence?}
\end{frame}

\begin{frame}[t]{E3: three models respond to the same numerical evidence}

\checkpoint{E3}{Actual public responses for the fixed first test scenario.}
\begin{block}{llama3.1:8b; 1.867 s}\fontsize{7.3}{8.3}\selectfont \textbf{INCREASE\_PV\_0025}. The incumbent candidate has a score of 0.031103819352072536, which is relatively low. Increasing active participation by 0.0025 may help reduce curtailment and improve the overall performance. This is supported by the operation's description, which states that it increases active participation to reduce curtailment.
\end{block}
\begin{block}{llama3.2:3b; 0.905 s}\fontsize{7.3}{8.3}\selectfont \textbf{RELAX\_Q\_10\_PERCENT}. Reducing reactive effort magnitude by ten percent may help improve the overall system performance and reduce curtailment, as evidenced by the current step J calculation which includes a term of 0.05*b\textasciicircum{}2, indicating a positive impact of reducing q on the objective function.
\end{block}
\begin{block}{qwen2.5:7b-instruct; 1.844 s}\fontsize{7.3}{8.3}\selectfont \textbf{RELAX\_Q\_10\_PERCENT}. The current incumbent solution has a lower objective score compared to the SAC candidate, indicating it is performing better. Relaxing the reactive effort by 10\% may help improve voltage stability without significantly degrading active power performance, as evidenced by the SAC candidate's higher score.
\end{block}

\slidetransition{How are different phrasings grouped into physical actions?}
\sources{Executed voltage pilot: \texttt{LLM\_Response\_Selection.ipynb}; frozen protocols and actual response journals: \texttt{results/ensemble\_voltage\_pilot/}. Numerical evidence and final checks are retained.}
\note{Optional historical/detail reference, outside the current 180-minute route. Original notebook IDs and experiment scopes are preserved. Keep the reasons verbatim and ask students which statements can be checked against available evidence. These are public model outputs, not hidden reasoning. Test 01 is a fixed illustration; the notebook exposes every response and reports aggregate outcomes. Identical prompts and distinct model configurations do not establish independent errors. Spoken bridge: How are different phrasings grouped into physical actions?}
\end{frame}

\begin{frame}[t]{E3: canonical actions make the vote reproducible}

\checkpoint{E3}{Inspect the actual operation, votes and physical controls.}
\cellcode{E3}{1}{7}
\small\medskip
\textbf{Canonical key:} tool name plus finite structured control parameters.\par\smallskip
\textbf{Evidence limit:} Qwen mentions voltage stability; these tests establish only steady-state quantities.\par\smallskip
\textbf{Majority tie rule:} smallest canonical SHA-256 ID, fixed before evaluation.\par\smallskip
\textbf{Baseline cost:} fixed pool response, majority and learned selection all pay for the same three calls.

\slidetransition{When does the selector receive its numerical reward?}
\sources{Executed voltage pilot: \texttt{LLM\_Response\_Selection.ipynb}; frozen protocols and actual response journals: \texttt{results/ensemble\_voltage\_pilot/}. Numerical evidence and final checks are retained.}
\note{Optional historical/detail reference, outside the current 180-minute route. Original notebook IDs and experiment scopes are preserved. Run the whole E3 cell, not just the excerpt. Canonicalization rejects malformed replies and duplicate model identities. The SHA tie-break is deterministic and has no physical preference. A fixed response from the common pool is a controlled selection comparator, not a single-model deployment runtime. Full raw call payloads remain in the journal. Spoken bridge: When does the selector receive its numerical reward?}
\end{frame}

\begin{frame}[t]{E4: the reward is revealed only after a response is selected}

\checkpoint{E4}{Reproduce the first actual training update.}
\cellcode{E4}{4}{12}
\fontsize{8.1}{9.5}\selectfont
\[
r=\begin{cases}
\min(1,\max(0,J_{\rm inc}-J_{\rm trial})/0.01)-0.01, & \text{trial passes},\\
-1.01, & \text{trial rejected}.
\end{cases}
\]
$J_{\rm inc}$: incumbent cost; $J_{\rm trial}$: selected-trial cost. The fixed trial penalty is not seconds. Numerical acceptance remains a separate hard check.

\slidetransition{How do training records differ from frozen evaluation?}
\sources{Executed voltage pilot: \texttt{LLM\_Response\_Selection.ipynb}; frozen protocols and actual response journals: \texttt{results/ensemble\_voltage\_pilot/}. Numerical evidence and final checks are retained. Bandit method: \href{https://arxiv.org/abs/1003.0146}{Li et al. (2010)}.}
\note{Optional historical/detail reference, outside the current 180-minute route. Original notebook IDs and experiment scopes are preserved. This is executable notebook code, and its first operation/reward is checked against the recorded training trace. The outcomes exist in a separate file for audit, but the operation is selected before that file is opened. No test outcome enters training, and the independent AC audit is excluded from the MV reward. All response-generation cost has already been incurred, so the reward does not pretend that selecting one response saves the other calls. Spoken bridge: How do training records differ from frozen evaluation?}
\end{frame}

\begin{frame}[t]{E4: actual learning records and frozen evaluation}

\checkpoint{E4}{Five predeclared training orders over the same 24 contexts.}
\begin{center}\includegraphics[width=.96\linewidth,height=.44\textheight,keepaspectratio]{\figfile{selector_E4.png}}\end{center}
\fontsize{8.3}{9.8}\selectfont
Each checkpoint has 24 reward updates. Validation and test update $A$ and $b$ zero times. All five orders yield the same held-out MV actions and objectives here.
\par\smallskip This small regularized linear fit demonstrates learning mechanics. The trace does not establish convergence or generalization across feeders.

\slidetransition{Does the learned selector beat voting on held-out cases?}
\sources{Executed voltage pilot: \texttt{LLM\_Response\_Selection.ipynb}; frozen protocols and actual response journals: \texttt{results/ensemble\_voltage\_pilot/}. Numerical evidence and final checks are retained.}
\note{Optional historical/detail reference, outside the current 180-minute route. Original notebook IDs and experiment scopes are preserved. SAC training and response-selection training have separate checkpoints. These curves show the bandit, not the 12,000-transition SAC training experiment. Training rewards can be negative even for a feasible trial because the improvement is small relative to the fixed trial penalty. The uncertainty bonus supports exploration during training; frozen evaluation uses predicted reward only. Spoken bridge: Does the learned selector beat voting on held-out cases?}
\end{frame}

\begin{frame}[t]{E5: learned selection ties voting on the held-out cases}

\checkpoint{E5}{Same three responses, one selected trial, fresh final check.}
\begin{center}\includegraphics[width=.99\linewidth,height=.54\textheight,keepaspectratio]{\figfile{selector_E5.png}}\end{center}
\fontsize{8.3}{9.8}\selectfont
MV test: learned selection and voting tie in all 12 cases. Mean $J=0.03398479$ versus numerical incumbent $0.03400656$; the reduction is small.\par
Every retained MV action fails the separate full-feeder AC audit.

\slidetransition{Can a rejected trial still end in a feasible recommendation?}
\sources{Executed voltage pilot: \texttt{LLM\_Response\_Selection.ipynb}; frozen protocols and actual response journals: \texttt{results/ensemble\_voltage\_pilot/}. Numerical evidence and final checks are retained.}
\note{Optional historical/detail reference, outside the current 180-minute route. Original notebook IDs and experiment scopes are preserved. This is the central selector result. The interval against voting collapses to zero because every paired outcome is identical. The graph gives all-responses extra trial budget and labels it accordingly. A small improvement over a particular fixed operation is not evidence that learning or LLM generation is necessary. E5 also displays a post-hoc diagnostic of every permitted fixed numerical rule; always relaxing Q matches the selected pattern without language calls. The twelve contexts are the statistical units. Spoken bridge: Can a rejected trial still end in a feasible recommendation?}
\end{frame}

\begin{frame}[t]{E5: trial rejection and final quality are different outcomes}

\checkpoint{E5--E6}{Compare each method within its declared model scope.}
\fontsize{8.7}{10.2}\selectfont
\begin{tabular}{@{}llrrr@{}}\toprule Backend & Method & Mean final $J$ & Trial rejects & Final AC PASS \\ \midrule
MV & Fixed numerical & 0.03400656 & 0/12 & 0/12 \\
MV & Majority vote & 0.03398479 & 0/12 & 0/12 \\
MV & Learned selector & 0.03398479 & 0/12 & 0/12 \\
AC & Fixed numerical & 0.47112138 & 11/12 & 12/12 \\
AC & Majority vote & 0.47138488 & 0/12 & 12/12 \\
AC & Learned selector & 0.47138488 & 0/12 & 12/12 \\

\bottomrule\end{tabular}\par\medskip
\textbf{AC transfer:} the learned selector avoids rejected trials but has a worse mean objective than fixed numerical selection. One larger missed improvement matters.\par\medskip
\textbf{Same gate:} a rejected refinement retains the incumbent, which can still pass its final check.\par\medskip
The final AC PASS column uses the broader-model audit for MV and the selected AC gate for OpenDSS. Different model scopes also change the voltage term in $J$.

\slidetransition{Which constraints become visible in the broader model?}
\sources{Executed voltage pilot: \texttt{LLM\_Response\_Selection.ipynb}; frozen protocols and actual response journals: \texttt{results/ensemble\_voltage\_pilot/}. Numerical evidence and final checks are retained.}
\note{Optional historical/detail reference, outside the current 180-minute route. Original notebook IDs and experiment scopes are preserved. Keep the objective and rejection metrics separate. The fixed increase operation is often rejected under AC constraints but in one context finds a larger improvement. Majority and learned selection favor the smaller Q refinement and are identical in this transfer set. Zero rejected trials therefore does not imply the best objective. All numbers are conditional on the declared method and twelve scenarios; final full-feeder PASS is not a field guarantee. Spoken bridge: Which constraints become visible in the broader model?}
\end{frame}

\begin{frame}[t]{E6: the model determines which constraints are visible}

\checkpoint{E6}{The same 12 physical contexts under both planning models.}
\begin{center}\includegraphics[width=.94\linewidth,height=.50\textheight,keepaspectratio]{\figfile{selector_E6.png}}\end{center}
\fontsize{8.5}{10}\selectfont
Learned workflow: \textbf{0/12} retained MV actions pass the full-feeder AC audit. With OpenDSS planning and verification, \textbf{12/12} final actions pass.\par
MV aggregation omits internal LV/transformer constraints. This comparison changes constraint coverage; it is not an accuracy-only test of two equivalent models.

\slidetransition{What inference and solver time does each method consume?}
\sources{Executed voltage pilot: \texttt{LLM\_Response\_Selection.ipynb}; frozen protocols and actual response journals: \texttt{results/ensemble\_voltage\_pilot/}. Numerical evidence and final checks are retained.}
\note{Optional historical/detail reference, outside the current 180-minute route. Original notebook IDs and experiment scopes are preserved. The plotted controls are the initial MV incumbent and separately optimized OpenDSS MPC plan, while the sentence reports final learned-workflow gates. The notebook exposes both. Avoid claiming the learned selector caused the AC benefit: the numerical model and acceptance contract changed. The selector trained on MV scenarios is frozen for AC transfer, so this is not separately trained in-domain learning on AC. Operator review must show both scope and failed broader-model checks. Spoken bridge: What inference and solver time does each method consume?}
\end{frame}

\begin{frame}[t]{E7: actual inference time belongs in the comparison}

\checkpoint{E7}{Recorded calls, component accounting and cached replay differ.}
\begin{center}\includegraphics[width=.97\linewidth,height=.50\textheight,keepaspectratio]{\figfile{selector_E7.png}}\end{center}
\fontsize{8.2}{9.8}\selectfont
168 actual calls, zero call/schema failures. MV mean accounted serial time: learned $8.119$ s; fixed numerical $3.327$ s.\par
All three responses are already generated when selection begins. The learner cannot save their inference time. AC transfer discloses reused MPC planning separately.

\slidetransition{What should the operator see after this comparison?}
\sources{Executed voltage pilot: \texttt{LLM\_Response\_Selection.ipynb}; frozen protocols and actual response journals: \texttt{results/ensemble\_voltage\_pilot/}. Numerical evidence and final checks are retained.}
\note{Optional historical/detail reference, outside the current 180-minute route. Original notebook IDs and experiment scopes are preserved. The plot shows actual local call latency, including model loading and residency effects. Calls were collected in model batches. Total policy durations are reconstructed sums of observed components, not direct online end-to-end stopwatch measurements. Fixed numerical generates no language pool. Cached notebook execution makes no new model calls and must not be reported as inference speed. Do not compare AC component totals excluding reused planning with MV totals including planning. Spoken bridge: What should the operator see after this comparison?}
\end{frame}

\begin{frame}[t]{E8: the operator sees the recommendation and its scope}

\checkpoint{E8}{Inspect the actual first test scenario and all five stage outputs.}
\fontsize{8.6}{10.2}\selectfont
\begin{tabular}{@{}p{.30\linewidth}p{.30\linewidth}p{.31\linewidth}@{}}
\toprule Operator field & MV planning workflow & OpenDSS transfer workflow \\
\midrule
Selected-model final check & PASS & PASS \\
Full-feeder check & FAIL & PASS \\
Advisory status & HOLD for broader-model review & Ready for human review \\
Displayed evidence & Control, model scope, failed AC constraints & Control, final constraints and candidate identity \\
Human acknowledgment & Pending & Pending \\
Equipment dispatch & None & None \\
\bottomrule\end{tabular}\par\medskip
\textbf{Five returned stage records:} context; MPC/SAC controls; fresh simulation metrics; gate result; selected response operation and measured gain.\par\medskip
The existing P dashboard remains a separate historical readback exercise. E records do not silently replace its recommendations.

\slidetransition{Can a feasible incumbent still benefit from one bounded trial?}
\sources{Executed voltage pilot: \texttt{LLM\_Response\_Selection.ipynb}; frozen protocols and actual response journals: \texttt{results/ensemble\_voltage\_pilot/}. Numerical evidence and final checks are retained.}
\note{Optional historical/detail reference, outside the current 180-minute route. Original notebook IDs and experiment scopes are preserved. The cell creates a compact operator report with exact candidate and evidence hashes. A full-feeder recommendation is held when the MV action fails the separate AC audit. The accepted AC action still requires human review and an authorized operating process. An acknowledgment is not generated by the experiment. The notebook also prints each stage method and actual returned output, making the abstraction traceable to recorded code and evidence. Spoken bridge: Can a feasible incumbent still benefit from one bounded trial?}
\end{frame}

\begin{frame}[t]{A passed candidate can still be improved}
\begin{center}
\begin{tikzpicture}[box/.style={draw=forest,rounded corners,align=center,text width=2.65cm,minimum height=1.05cm,font=\small},flow/.style={->,>=stealth,thick,forest}]
\node[box] (model) at (0,0) {\textbf{Simulation}\\Predict the outcome};
\node[box] (gate) at (3.8,0) {\textbf{Verification}\\Check limits and policy};
\node[box,fill=mint!20] (pool) at (7.6,0) {\textbf{Feasible shortlist}\\Retain cost and evidence};
\node[box,fill=mint!35] (opt) at (3.8,-2.0) {\textbf{LLM Optimizer}\\Repair or refine};
\node[box] (rank) at (11.0,-2.0) {\textbf{Rank + recheck}\\Select verified lowest cost};
\draw[flow] (model.east)--(gate.west);
\draw[flow] (gate.east)--node[above,font=\scriptsize]{PASS}(pool.west);
\draw[flow,orange] (gate.south)--node[left,font=\scriptsize]{REJECT}(opt.north);
\draw[flow] (pool.south)|-node[below,pos=.74,font=\scriptsize]{Try an improvement}(opt.east);
\draw[flow,orange] (opt.west)-|node[below,pos=.28,font=\scriptsize]{Changed candidate}(model.south);
\draw[flow] (pool.east)-|node[right,pos=.72,font=\scriptsize]{Stop / budget}(rank.north);
\end{tikzpicture}
\end{center}
\plainidea{Only accepted candidates enter the ranking. Every changed action returns to simulation and verification.}
\slidetransition{How do we enforce an equal maximum trial budget?}
\sources{SmartGridComm 2026 tutorial proposal, p. 2: repair, refinement and simulation-informed ranking. Executed state machine: \texttt{src/proposal\_workflow.py}.}
\note{Optional historical/detail reference, outside the current 180-minute route. Original notebook IDs and experiment scopes are preserved. Use this after the five-module overview. Awareness and Planning have already produced the initial candidates. Simulation and Verification create candidate-specific evidence. A rejected candidate supplies a reason for repair; an accepted candidate may still have poor objective value and invite refinement. The fifth proposal module handles both branches. Under contract revision 1.1, the first bounded refinement is required when valid tools are available. A later iteration may stop through KEEP_INCUMBENT, an exhausted operation set or the maximum budget. This scheduled first refinement is a shared protocol rule, not a freely selected agent behavior. Call or tool failures still lead to REVIEW. A valid shortlist is ranked, and the chosen action is simulated and verified once more before recommendation. If the process requests review or the final check fails, no recommendation is issued. The older eight named helper responsibilities organize execution and explanation; they do not replace these five core proposal modules. Spoken bridge: How do we enforce an equal maximum trial budget?}
\end{frame}

\begin{frame}[t]{Optimize over tested candidates under the same maximum budget}
\[
\mathcal F_B=\{u\in\mathcal C_B:\ \mathrm{solved}(u),\ g_j(x(u),u)\leq0\ \forall j,\ \mathrm{policy}(u)\},
\qquad
u^\star\in\underset{u\in\mathcal F_B}{\arg\min}\ J(x(u),u).
\]
\small
\begin{tabular}{@{}>{\raggedright\arraybackslash}p{.18\linewidth}>{\raggedright\arraybackslash}p{.74\linewidth}@{}}
$\mathcal C_B,\mathcal F_B$ & Tested search candidates and the subset passing all declared checks. \\
$u,x(u),J$ & Action, simulated state and case-specific cost to minimize. \\
$g_j,j,\mathrm{policy}$ & Engineering-limit residual, check index and permitted-action rules. \\
$B,u^\star$ & Maximum new refinement evaluations; best tested feasible action. \\
\end{tabular}
\medskip
\begin{center}
\begin{tikzpicture}[box/.style={draw=forest,rounded corners,align=center,text width=3.65cm,minimum height=.9cm,font=\small},flow/.style={->,>=stealth,thick,forest}]
\node[box] (seeds) at (0,0) {\textbf{2 shared seeds}\\Same tools\\and constraints};
\node[box,fill=mint!25] (refine) at (4.65,0) {\textbf{At most $B=2$ trials}\\First refinement required*};
\node[box] (final) at (9.3,0) {\textbf{1 final recheck}\\If an action is recommended};
\draw[flow] (seeds.east)--(refine.west); \draw[flow] (refine.east)--(final.west);
\end{tikzpicture}
\end{center}
\par\scriptsize *Required when a valid refinement tool is available; later search may stop.\\
Same maximum budget does not imply equal actual evaluations. Report actual counts and elapsed time.\\
Internal MPC solves and the initial state snapshot are additional work, counted separately.
\slidetransition{Which claims belong to the public reason and which to the solver?}
\sources{Finite evaluated-set selection and matched protocol: \texttt{src/proposal\_workflow.py}, \texttt{src/proposal\_case\_tools.py}. Proposal, p. 2: best-of-$N$ and verifier-guided refinement.}
\note{Optional historical/detail reference, outside the current 180-minute route. Original notebook IDs and experiment scopes are preserved. This is a common finite-search formulation, not a guarantee of the global continuous or mixed-integer optimum. The numerical adapters define J, the model and the constraints separately for each physical task. Contract revision 1.1 allows at most two refinement evaluations after two shared seed candidates and requires the first refinement when valid tools are available. Later search may stop early, so actual counts can differ. The fixed and language-guided routes share this rule; the language model chooses among permitted operations rather than freely choosing whether to perform the required first refinement. A failed call or numerical tool can stop with REVIEW. The retained initial contract pilot and revised run must not be pooled as a repeated benchmark. A final independent recheck is outside the search set and is recorded separately; it duplicates the selected physical action and does not add a new optimization candidate. Awareness and MPC preparation can call the electrical model internally, so do not label the four-candidate search cap as a four-power-flow cap. An empty feasible set or held review produces REVIEW. Compare accepted objective values, budget use, failures and full workflow cost, not ticks alone. Spoken bridge: Which claims belong to the public reason and which to the solver?}
\end{frame}

\begin{frame}[t]{The public Llama reason and the numerical evidence have different jobs}
\begin{columns}[T,onlytextwidth]
\begin{column}{.47\textwidth}
\begin{block}{Llama response: interpret the evidence}
\textbf{choice}: one permitted operation.\par\smallskip
\textbf{reason}: a short returned explanation.\par\smallskip
\textbf{ranking}: verified candidate IDs only.
\end{block}
\end{column}
\begin{column}{.49\textwidth}
\begin{block}{Numerical record: check the claim}
Candidate and operating-contract identity.\par\smallskip
Fresh metrics, violations and cost.\par\smallskip
Independent ordering by numerical cost.
\end{block}
\end{column}
\end{columns}
\medskip
\begin{center}
\begin{tikzpicture}[box/.style={draw=forest,rounded corners,align=center,text width=3.5cm,minimum height=.92cm,font=\small},flow/.style={->,>=stealth,thick,forest}]
\node[box] (schema) at (0,0) {Validate the choice\\and every candidate ID};
\node[box,fill=mint!25] (audit) at (4.55,0) {Check the order\\against numerical costs};
\node[box] (choose) at (9.1,0) {Select numerical best\\then recheck};
\draw[flow] (schema.east)--(audit.west); \draw[flow] (audit.east)--(choose.west);
\end{tikzpicture}
\end{center}
\par\small A fluent explanation is a recorded claim. Its supporting measurements and checks remain visible.
\slidetransition{Where is the bounded decision loop implemented?}
\sources{Tool-observation feedback: \href{https://arxiv.org/abs/2210.03629}{Yao et al., ReAct (2023)}. Response validation and ranking guard: \texttt{src/proposal\_workflow.py}.}
\note{Optional historical/detail reference, outside the current 180-minute route. Original notebook IDs and experiment scopes are preserved. No invented model quotation is shown on this concept slide. Later slides display actual returned reason strings and their call IDs. These are public explanations generated by Llama, not hidden internal chain-of-thought and not a physical proof. The shared contract schedules a required first refinement when valid tools are available; Llama chooses its allowed operation, not the existence of every stage. The ranking schema permits every feasible ID exactly once. The implementation records whether Llama's ordering agrees with cost, but numerical minimum cost determines the selected action; language cannot promote a rejected candidate or override the objective. Wrong structure or IDs force review. A cost-order disagreement remains visible for the operator and evaluation rather than being silently presented as correct language reasoning. Latency, token counts, model settings and errors belong to the same trace record. Spoken bridge: Where is the bounded decision loop implemented?}
\end{frame}

\begin{frame}[t]{Initial supervisor code 1/9: a bounded loop controls what can happen next}

\implcode{continuous_workflow}{217}{223}
\footnotesize
\begin{tabular}{@{}>{\raggedright\arraybackslash}p{.22\linewidth}>{\raggedright\arraybackslash}p{.70\linewidth}@{}}
\textbf{Attempt 1} & Validate the shared SAC proposal and request a fresh simulation. \\
\textbf{Attempt 2} & After rejection, check the eligible MPC first action or stop for review. \\
\textbf{Terminal result} & Return a recommendation within the selected model's scope, or return \texttt{REVIEW}. There is no third attempt or equipment command. \\
\end{tabular}
\par\smallskip Both supervisors use the same numerical route and acceptance limits. Llama chooses only a permitted next step.
\slidetransition{How does the operator request become a numerical contract?}
\sources{Executable transition rules: \texttt{src/continuous\_workflow.py}, \texttt{run\_continuous\_workflow}, including lines 252--265. Historical continuous-supervisor protocol 1.1; the new five-stage loop is a separate experiment.}
\note{Optional historical/detail reference, outside the current 180-minute route. Original notebook IDs and experiment scopes are preserved. Optional implementation detail. Begin with the literal range(2), rather than another abstract architecture drawing. The initial planner is allowed VERIFY_SAC_FIRST or REVIEW. After a failed first check, a validated MPC plan can be offered once. A planner REVIEW is terminal for candidate evaluation. Execution here means a numerical proposal, not field dispatch. Ask the audience to locate the two restrictions that prevent arbitrary tool loops: the finite attempt count and the expected tool name. Spoken bridge: How does the operator request become a numerical contract?}
\end{frame}

\begin{frame}[t]{Initial supervisor code 2/9: the request becomes a numerical contract}

\implcode{continuous_workflow}{49}{56}
\footnotesize
\textbf{Units:} voltage is per unit; current allowance is percent of branch rating. For inverter $g$, $S_g$ is kVA, $p_g$ is kW, $q_g$ is kvar, and $a,b$ are dimensionless fleet participation factors.
\par\smallskip\textbf{Authority:} the validated case supplies demand and limits. The language reply has no field for replacing them.
\par\smallskip\textbf{Notebook W2:} inspect this contract before a run. Correct misunderstood operator intent first.
\slidetransition{How do several roles share one local model service?}
\sources{Exact contract: \texttt{src/continuous\_workflow.py}. Case-domain validation and action bounds: \texttt{src/continuous\_control.py}, \texttt{validate\_case} and \texttt{\_control}. Notebook W2.}
\note{Optional historical/detail reference, outside the current 180-minute route. Original notebook IDs and experiment scopes are preserved. Optional. In the current P3 case, rating_fraction is 0.8, so the returned allowance is 80 percent, not an 80 percent reserve. The participation action is shared across the inverter fleet and is not an independently optimized command at every inverter. Contract validation checks form and declared numerical ranges. It does not establish that a request was interpreted correctly or that the sensors are authenticated. This is why readback remains necessary. Spoken bridge: How do several roles share one local model service?}
\end{frame}

\begin{frame}[t]{Initial supervisor code 3/9: named roles share one local Llama service}

\implcode{continuous_workflow}{69}{72}
\implcode{continuous_workflow}{190}{194}
\footnotesize
\textbf{Changes by role:} Planner sees the proposal and contingency plan. Recovery sees failed checks. Explanation sees the metric card. Each receives a restricted menu.
\par\smallskip\textbf{Shared:} one local \texttt{llama3.1:8b} backbone and the numerical tools. Sequential role prompts do not provide independent votes or independent safety assessment.
\par\smallskip The fixed route uses the same interface, so every recorded choice identifies its source.
\slidetransition{How is the language choice restricted to permitted tools?}
\sources{Actual client and role dispatch: \texttt{src/continuous\_workflow.py}. Role/tool interaction background: Yao et al., ReAct (ICLR 2023), \url{https://arxiv.org/abs/2210.03629}.}
\note{Optional historical/detail reference, outside the current 180-minute route. Original notebook IDs and experiment scopes are preserved. Optional. Explain that an agent architecture can impose useful separation of responsibilities without creating statistically independent judgments. This implementation makes no independence claim. One additional role can add another model call and another failure opportunity. The scientific question is whether the resulting interpretation or routing improves a measured task outcome after accounting for errors and cost. Spoken bridge: How is the language choice restricted to permitted tools?}
\end{frame}

\begin{frame}[t]{Initial supervisor code 4/9: the language tool chooses from a restricted menu}

\implcode{continuous_workflow}{75}{78}
\implcode{continuous_workflow}{81}{82}
\footnotesize
\begin{tabular}{@{}>{\raggedright\arraybackslash}p{.20\linewidth}>{\raggedright\arraybackslash}p{.72\linewidth}@{}}
\textbf{Planner menu} & \texttt{VERIFY\_SAC\_FIRST} or \texttt{REVIEW}. Requesting a check does not approve the proposal. \\
\textbf{Recovery menu} & \texttt{USE\_MPC\_PLAN} or \texttt{REVIEW}. The first choice still requires a fresh numerical check. \\
\textbf{Returned reason} & A public explanation of the selected choice, not numerical evidence or a correctness guarantee. \\
\end{tabular}
\slidetransition{Which checks follow valid JSON?}
\sources{Exact schema and system-prompt excerpt: \texttt{src/continuous\_workflow.py}. Structured-response API: \url{https://docs.ollama.com/api/chat}.}
\note{Optional historical/detail reference, outside the current 180-minute route. Original notebook IDs and experiment scopes are preserved. Optional. These excerpts are two separate portions of the actual ask method, not a reconstructed prompt. The call also includes definitions of the permitted actions from CHOICE_DEFINITIONS. The prompt states the policy, while program logic enforces the available routes and numerical gate. A prompt alone cannot establish electrical feasibility. The user-visible reason is stored model output, not a reconstruction of hidden reasoning. Spoken bridge: Which checks follow valid JSON?}
\end{frame}

\begin{frame}[t]{Initial supervisor code 5/9: valid JSON is only the first acceptance test}

\implcode{continuous_workflow}{99}{108}
\footnotesize
\textbf{Established:} two fields, a permitted choice and a nonempty reason. Service, decoding or schema failure produces local \texttt{REVIEW}.
\par\smallskip\textbf{Still unverified:} request semantics, measurement trust, optimality and voltage/current feasibility.
\par\smallskip\textbf{Provenance:} \texttt{actual\_local\_llama} is a parsed reply; \texttt{local\_failure\_fallback} is locally generated text. Neither proves physical safety.
\slidetransition{How does a passing result retain its model scope?}
\sources{Exact parser, fallback and client timer: \texttt{src/continuous\_workflow.py}. Response-source labels are returned at lines 112--114. API response structure: \url{https://docs.ollama.com/api/chat}.}
\note{Optional historical/detail reference, outside the current 180-minute route. Original notebook IDs and experiment scopes are preserved. Optional. The exception path is part of the workflow, not an LLM response. It must not be reported as successful model reasoning. Give the concrete counterexample from the intake study: a perfectly valid structured reply can misread a current reserve. A physical solver operating under that incorrect limit cannot recover the original human intent. The current numerical workflow receives an already approved structured case. Spoken bridge: How does a passing result retain its model scope?}
\end{frame}

\begin{frame}[t]{Initial supervisor code 6/9: a passing result always carries its model scope}

\implcode{continuous_workflow}{273}{278}
\footnotesize
\begin{tabular}{@{}>{\raggedright\arraybackslash}p{.22\linewidth}>{\raggedright\arraybackslash}p{.70\linewidth}@{}}
\textbf{OpenDSS} & Full supplied unbalanced MV/LV feeder. The experiment does not model protection coordination or transients. \\
\textbf{LinDistFlow3} & Coupled MV phase network with idealized LV aggregation. Internal LV voltage/current constraints are omitted. \\
\textbf{Both} & Fresh selected-model simulation, not field evidence or coverage of omitted equipment. \\
\end{tabular}
\par\smallskip\texttt{MODEL\_RECOMMEND} is scoped to the selected model. Different constraint coverage can produce different actions.
\slidetransition{What can a fresh solve reject that prose cannot detect?}
\sources{Exact returned scope: \texttt{src/continuous\_workflow.py}. Model assumptions: \texttt{src/lindistflow3.py} and \texttt{appendix/Stage\_Methods\_and\_Models.md}. Current recorded comparison: \texttt{results/continuous\_workflow\_comparison.json}.}
\note{Optional historical/detail reference, outside the current 180-minute route. Original notebook IDs and experiment scopes are preserved. Optional. Tie this source excerpt to the current P3 example: the MV model accepts the common SAC action, while full AC requires an MPC fallback. This does not establish that the linear MV prediction is inaccurate on the MV buses. It demonstrates why a result must expose its coverage. Plan-level screening also checks that the stored plan backend matches the selected backend. Fresh checking reuses the selected physical model, so it is not independent model validation. Spoken bridge: What can a fresh solve reject that prose cannot detect?}
\end{frame}

\begin{frame}[t]{Initial supervisor code 7/9: a fresh solve can reject a plausible proposal}

\implcode{continuous_workflow}{230}{237}
\footnotesize
\textbf{Inputs:} proposed $[a,b]$, first forecast demand multiplier and approved current allowance. The backend returns fresh action-specific evidence.
\par\smallskip\textbf{Gate:} finite, identified voltage/current arrays, convergence, scope and dispatch tracking. Missing or invalid evidence rejects.
\par\smallskip\textbf{Recovery:} exceptions and mismatches reject. An eligible MPC fallback passes through this same solve and gate.
\slidetransition{How does the code preserve the operator review boundary?}
\sources{Exact fresh-evaluation path: \texttt{src/continuous\_workflow.py}. Deterministic checks: \texttt{src/continuous\_control.py}, \texttt{verify\_control}. Here $a,b$ are dimensionless inverter participation factors.}
\note{Optional historical/detail reference, outside the current 180-minute route. Original notebook IDs and experiment scopes are preserved. Optional. Emphasize the difference between the model's own accepted flag and the deterministic gate that reconstructs the checks from the returned arrays. The checker requires nonempty finite data, identities, fresh convergence, valid scope, the requested case, voltage and current bounds, and actual versus commanded P/Q tracking. This is a first-action check. The MPC plan has separate horizon evidence, and neither proves stability or robustness to unknown disturbances. Spoken bridge: How does the code preserve the operator review boundary?}
\end{frame}

\begin{frame}[t]{Initial supervisor code 8/9: the final authority stays with the operator}

\implcode{continuous_workflow}{61}{64}
\implcode{continuous_workflow}{288}{291}
\footnotesize
\textbf{Read back:} proposed P/Q, participation factors, current allowance and scope. Review measurement, protection and implementation questions.
\par\smallskip\textbf{Record identity:} the hash binds the decision fields to this record. Local acknowledgment neither dispatches equipment nor overrides rejection.
\par\smallskip\textbf{Benefit to test:} consistent evidence for human review. Operator workload and live ADMS deployment are not evaluated.
\slidetransition{Which saved call produced the displayed response?}
\sources{Exact local-review contract: \texttt{src/continuous\_workflow.py}. Human/automation allocation background: Parasuraman, Sheridan and Wickens (2000), \url{https://doi.org/10.1109/3468.844354}. Dashboard checkpoint F2.}
\note{Optional historical/detail reference, outside the current 180-minute route. Original notebook IDs and experiment scopes are preserved. Optional. The operator remains a person, not another Llama role. The final dashboard is local and has no SCADA dispatch API. Hashing supports record identity but does not authenticate the human, sensors, model or source files. A future operational integration would need separate authority and assurance work. Ask the audience which displayed quantity would reveal a mistaken interpretation of a 20 percent current reserve. Spoken bridge: Which saved call produced the displayed response?}
\end{frame}

\begin{frame}[t]{Initial supervisor code 9/9: match the displayed response to its recorded call}

\checkpoint{T1}{Open the read-only trace notebook and load the supplied evidence.}
\implcode{continuous_workflow}{108}{114}
\footnotesize
\textbf{Stage record:} the role, numerical method, status, response source and returned content travel together. The journal retains the actual request context and raw local-model response.
\par\smallskip\textbf{Audit rule:} the reader requires the same call ID, choice, reason and recorded client duration. It rejects missing or nonunique matches instead of selecting a convenient journal.
\par\smallskip T2--T6 expose the voltage, MILP, h06 intake and cyber traces. T7--T9 then inspect timing fields and source hashes. Reading these files does not invoke a solver or Llama.
\slidetransition{Which two clocks appear in the recorded trace?}
\sources{Exact stage schema, journal and client timer: \texttt{src/continuous\_workflow.py}. Matching reader: \texttt{src/trace\_latency.py}. Companion: \texttt{Agent\_Trace\_and\_Latency.ipynb}, T1--T9.}
\note{Optional historical/detail reference, outside the current 180-minute route. Original notebook IDs and experiment scopes are preserved. Optional. This is the bridge from architecture to inspectable experiments: every claim on the later case slides should point to a numerical return or an actual recorded language response. The call ID is local to its journal, so the source path and unique matching rule also matter. Deterministic stages do not acquire invented LLM reasons or invented per-stage timings. The notebooks are evidence readers unless the main experiment explicitly launches a live call. Spoken bridge: Which two clocks appear in the recorded trace?}
\end{frame}

\begin{frame}[t]{Trace review 1/3: read token counts and two distinct clocks}

\checkpoint{T7}{Inspect the recorded calls and usage fields.}
\begin{columns}[T,onlytextwidth]
\begin{column}{.61\textwidth}
\includegraphics[width=\linewidth,height=.61\textheight,keepaspectratio]{\figfile{recorded_call_latency_slide.png}}
\end{column}
\begin{column}{.36\textwidth}
\footnotesize
\textbf{Client:} HTTP request, parsing and schema validation, timed in seconds.
\par\smallskip\textbf{Server:} \texttt{total\_duration} converted from nanoseconds. Do not add it to client time.
\par\smallskip\textbf{Tokens:} \texttt{prompt\_eval\_count} (input), \texttt{eval\_count} (output).
\par\smallskip\textbf{Example, current AC Planner:} 795 input tokens, 36 output tokens, 0.988 s client and 0.955 s server.
\par\smallskip Unmeasured stage times and causes stay unknown.
\end{column}
\end{columns}
\slidetransition{How does measured cost relate to the actual outcome?}
\sources{Matched raw journals read by \texttt{src/trace\_latency.py}, \texttt{\_call\_row}. API field units: \url{https://docs.ollama.com/api/chat}. Figure: three current OpenDSS calls; T7 includes all ten selected calls.}
\note{Optional historical/detail reference, outside the current 180-minute route. Original notebook IDs and experiment scopes are preserved. Optional. Explain that token throughput, if inspected in the table, is output count divided by eval_duration converted to seconds. It is a generation-rate measure, not total application throughput. For the current AC Planner this is about 74.3 generated tokens per second. Do not infer network overhead, queueing, hardware effects or cache causality from the difference of these aggregate intervals. The h06 client call is distinct from its slightly longer intake wrapper, and its downstream numerical runtime is a third interval. Spoken bridge: How does measured cost relate to the actual outcome?}
\end{frame}

\begin{frame}[t]{Trace review 2/3: compare cost with the recorded outcome}

\checkpoint{T8}{Compare fixed and agentic intervals within each recorded pair.}
\begin{columns}[T,onlytextwidth]
\begin{column}{.64\textwidth}
\includegraphics[width=\linewidth,height=.62\textheight,keepaspectratio]{\figfile{recorded_workflow_latency_slide.png}}
\end{column}
\begin{column}{.33\textwidth}
\footnotesize
\textbf{Fixed / agentic seconds}\par
AC: 0.346 / 3.110 s.\par
MV: 0.134 / 4.911 s.\par
Same action within each backend.
\par\smallskip\textbf{Historical cases}\par
MILP: 0.009859 / 6.10178 s.\par
Cyber: 0.002146 / 5.93465 s.\par
Checks pass; the agentic explanation returns \texttt{REVIEW}.
\par\smallskip\textbf{No speed gain.} Language calls add time and can add unnecessary handoffs.
\end{column}
\end{columns}
\par\footnotesize Current: precomputed controllers/training excluded. Historical: numerical tools included. One run per route; h06 at T5.
\slidetransition{What should the next controlled study measure?}
\sources{Records: \texttt{results/workshop\_run\_[backend].json}, \texttt{standard\_agent\_workflow.json}. Interval definitions: \texttt{src/trace\_latency.py}.}
\note{Optional historical/detail reference, outside the current 180-minute route. Original notebook IDs and experiment scopes are preserved. Optional. These are descriptive observations on one machine, not matched repeated statistical benchmarks. Compare fixed and agentic within each recorded pair, but do not pool protocols or rank backend speed from differently timed runs. The difference between total workflow time and summed client calls is only other work inside that interval, not a measured solver-stage timer. The positive h06 intake example demonstrates language coverage with errors elsewhere in the study; it does not turn these timing results into a speed advantage. Spoken bridge: What should the next controlled study measure?}
\end{frame}

\begin{frame}[t]{Trace review 3/3: audit this record; design the next study}

\checkpoint{T9}{Inspect the source hashes and evidence audit.}
\cellcode{T9}{1}{5}
\footnotesize
\begin{tabular}{@{}>{\raggedright\arraybackslash}p{.24\linewidth}>{\raggedright\arraybackslash}p{.68\linewidth}@{}}
\textbf{What this audit proves} & Matches specific saved bytes and replies. A hash detects change against a retained reference. \\
\textbf{What it cannot prove} & Truth, model adequacy, operator identity or general performance. Recomputed hashes do not authenticate sources. \\
\textbf{Next matched study} & Freeze cases, limits, models, prompts and code. Repeat both routes with declared seeds and randomized order. \\
\textbf{Measure together} & Correctness, acceptance, unnecessary review, failures and time. Separate cold/warm setup and measure operator outcomes. \\
\end{tabular}

\slidetransition{How are the three phases coupled in the linear approximation?}
\sources{Exact source hashing: \texttt{src/trace\_latency.py}. Raw-source hashes and unique journal matching are returned by \texttt{load\_recorded\_evidence}. Companion T9 records the files actually inspected.}
\note{Optional historical/detail reference, outside the current 180-minute route. Original notebook IDs and experiment scopes are preserved. Optional. There are eight source files in the current evidence collection, but the count is not hard-coded here so the notebook remains authoritative if packaging adds records. Store an external trusted reference if authenticity matters. Distinguish a file hash from the recommendation-field hash used by the dashboard. For a repeated experiment, predeclare handling of timeouts, failed schemas, incorrect READY replies and correct abstentions. Do not silently drop slow calls or failed tasks before reporting latency. This page proposes future methodology; it claims no new measurements. Spoken bridge: How are the three phases coupled in the linear approximation?}
\end{frame}

\begin{frame}[t]{Optional reference: coupled three-phase LinDistFlow}
\referencefont
\[
S_{ij}=s_j+\sum_{k:j\to k}S_{jk},\qquad
w_j=w_i-2\operatorname{Re}\{\overline D_\beta Z_{ij}D_\beta\overline S_{ij}\}.
\]
\referencefont
\begin{tabular}{@{}>{\raggedright\arraybackslash}p{.24\linewidth}>{\raggedright\arraybackslash}p{.69\linewidth}@{}}
\toprule Symbol / domain & Meaning and units \\
\midrule
$i,j,k$; $j\to k$ & Parent bus, child bus and downstream child; sum over the outgoing branches of bus $j$. \\
$S_{ij}\in\mathbb C^{n_\phi}$ & Branch phase powers $P_{ij}+\mathrm{i}Q_{ij}$ in VA (real parts W, imaginary parts var); $n_\phi$ is the branch's number of present phases. \\
$s_j\in\mathbb C^{n_\phi}$ & Net complex consumption at bus $j$, positive for consumption and negative for injection, in VA. \\
$w_i,w_j\in\mathbb R_+^{n_\phi}$ & Squared phase-to-neutral voltage magnitudes $|V|^2$, in V$^2$; the solver later converts magnitudes to per unit. \\
$Z_{ij}\in\mathbb C^{n_\phi\times n_\phi}$ & Phase impedance matrix in $\Omega$, including mutual impedances. \\
$D_\beta=\operatorname{diag}(\beta)$ & Nominal phase ratios $\beta=(1,e^{-\mathrm{i}2\pi/3},e^{\mathrm{i}2\pi/3})$ restricted to present phases; dimensionless. \\
$\overline{(\cdot)},\operatorname{Re},\mathrm{i}$ & Complex conjugation, componentwise real part, and imaginary unit $\mathrm{i}^2=-1$. \\
\bottomrule
\end{tabular}
\par\smallskip\textbf{Assumptions:} radial energized topology; nearly nominal phase-angle ratios; series-loss terms omitted in the balance. Balance terms use a common bus-phase ordering with absent phases zero-filled; voltage drop uses present branch phases. The export aggregates LV injections and omits internal LV constraints.
\slidetransition{How do objective terms and constraints define voltage support?}
\sources{Gan and Low (2014), Sec. V, \url{https://arxiv.org/abs/1406.3054}. Export and implementation: \texttt{src/lindistflow3.py}.}
\note{Optional historical/detail reference, outside the current 180-minute route. Original notebook IDs and experiment scopes are preserved. Optional, outside the three-hour timing. Explain dimensions before the multiplication: ohms times VA gives volts squared. The displayed equation uses physical units; multiplying kVA directly by ohms without converting would be wrong. The real-part expression retains phase coupling. Approximate reported losses are separate from the lossless balance. Model coverage, approximation error and solver convergence are different questions. Spoken bridge: How do objective terms and constraints define voltage support?}
\end{frame}

\begin{frame}[t]{Optional reference: voltage-control objective and constraints}
\referencefont
\[
\min_{x_k,u_k}\sum_{k,i,\phi}w_v(v_{i\phi,k}-v^{\rm ref})^2
+\sum_{k,g}\big[w_p(p^{\rm av}_{g,k}-p_{g,k})^2+w_q q_{g,k}^2\big]
\]
\[
F_m(x_k,u_k;\hat d_k)=0,\quad \underline v\leq v\leq\overline v,\quad
|I|\leq\overline I,\quad 0\leq p_g\leq p_g^{\rm av},\quad p_g^2+q_g^2\leq S_g^2.
\]
\referencefont
\begin{tabular}{@{}>{\raggedright\arraybackslash}p{.25\linewidth}>{\raggedright\arraybackslash}p{.68\linewidth}@{}}
$k,i,\phi,g$ & Time-step, bus, available-phase and PV-inverter indices. Sums cover the modeled horizon, bus phases and inverters. \\
$x_k\in\mathbb R^{n_x},u_k\in\mathbb R^{2n_g}$ & Real representation of electrical states and inverter P/Q controls; $n_x$ is the state dimension and $n_g$ the inverter count. \\
$F_m,\hat d_k$ & Selected model $m$'s power-flow residual and forecast active/reactive demand (kW/kvar). Zero residual imposes its electrical equations. \\
$v,v^{\rm ref},\underline v,\overline v$ & Voltage magnitude, target and bounds in pu; $I,\overline I$ are phase-current magnitude and allowed limit in A. Inequalities are componentwise. \\
$p_g,p_g^{\rm av},q_g,S_g$ & Active output/availability (kW), reactive output (kvar), apparent rating (kVA). \\
$w_v,w_p,w_q\geq0$ & Objective weights; their units or input normalization must make the summed objective consistent. \\
\end{tabular}
\par\smallskip\textbf{Executed snapshot reduction:} $p_g=aS_g$, $q_g=bS_g$ with dimensionless $a\in[0,1]$, $b\in[-.44,.44]$ and $a^2+b^2\leq1$. It minimizes the dimensionless cost $(1-a)^2+.05b^2+10\operatorname{mean}[(v-1.01)^2]$ across modeled voltage entries.
\slidetransition{What exactly is the implemented finite-horizon plan?}
\sources{Maharjan et al. (2021), \url{https://doi.org/10.1109/TSTE.2020.3001115}. Restricted snapshot: \texttt{src/standard\_formulations.py}.}
\note{Optional historical/detail reference, outside the current 180-minute route. Original notebook IDs and experiment scopes are preserved. Optional. The first equations describe a general problem; the final paragraph defines the actual two-variable reduction. Available P equals rated S only under the supplied high-irradiance overlay. A sequential local sensitivity solve, inscribed capability polygon and fresh checks do not provide global AC optimality. The general model can be extended to individual inverter decisions, but that larger optimization was not executed here. Spoken bridge: What exactly is the implemented finite-horizon plan?}
\end{frame}

\begin{frame}[t]{Optional reference: the executed continuous horizon plan}
\referencefont
\[
\min_{u_0,u_1,u_2}\sum_{k=0}^{2}\gamma^k\ell_k,\quad
\ell_k=(1-a_k)^2+.05b_k^2+\frac{10}{N_v}\sum_{r=1}^{N_v}(v_{r,k}-1.01)^2
+.05\|u_k-u_{k-1}\|_2^2.
\]
\referencefont
\begin{tabular}{@{}>{\raggedright\arraybackslash}p{.25\linewidth}>{\raggedright\arraybackslash}p{.68\linewidth}@{}}
\toprule Symbol / value & Meaning and implementation \\
\midrule
$H=3$, $k\in\{0,1,2\}$ & Three forecast steps. No physical step duration or dynamic energy state is specified in this quasistatic teaching model. \\
$u_k=(a_k,b_k)\in\mathbb R^2$ & Dimensionless fleet participation; $p_g=a_kS_g$ kW and $q_g=b_kS_g$ kvar for inverter rating $S_g$ in kVA. \\
$u_{-1}$ & Declared previous fleet command; the Euclidean squared norm penalizes changes from it or from the preceding planned step. \\
$v_{r,k}$; $N_v$ & Per-unit voltage magnitude at modeled bus-phase entry $r$; number of such entries in the selected model. \\
$\gamma=.98$, $\ell_k$ & Dimensionless discount and dimensionless engineering stage cost. Numeric coefficients are declared teaching weights, not fitted monetary costs. \\
$\rho\in(0,1]$ & Permitted nominal-current fraction: each predicted phase current must satisfy $|I|\leq\rho I^{\rm rated}$ (A). \\
\bottomrule
\end{tabular}
\par\smallskip\textbf{Algorithm:} rebuild local voltage/current sensitivities, solve a constrained QP step, and retain only a nonincreasing-cost horizon passing fresh model checks. The 32-sided inscribed capability polygon is conservative. The tool proposes $u_0$; this study inspects a horizon decision rather than an executed closed-loop rollout of MPC.
\slidetransition{How do policy outputs map to physical inverter controls?}
\sources{Rawlings, Mayne and Diehl, \emph{Model Predictive Control}, 2nd ed. (2017), \url{https://sites.engineering.ucsb.edu/~jbraw/mpc/}. Executed plan: \texttt{src/continuous\_control.py}.}
\note{Optional historical/detail reference, outside the current 180-minute route. Original notebook IDs and experiment scopes are preserved. Optional. Six real control variables are optimized. The three forecast loads and movement penalty couple the stages, but storage dynamics are absent. Fresh horizon feasibility is tool evidence; the common workflow gate checks the chosen first action. There is no recursive-feasibility, stability, uncertainty-robustness or global-optimality certificate. A small voltage-objective value on an MV-only model cannot certify LV service voltage. Spoken bridge: How do policy outputs map to physical inverter controls?}
\end{frame}

\begin{frame}[t]{Optional reference: SAC objective, action mapping and reward}
\setlength{\parskip}{0pt}
\referencefont
\setlength{\abovedisplayskip}{0pt}\setlength{\belowdisplayskip}{1pt}\setlength{\abovedisplayshortskip}{0pt}\setlength{\belowdisplayshortskip}{1pt}
\[
\max_\pi\;\mathbb E_\pi\!\left[\sum_{t=0}^{2}\gamma^t\{r_t+\alpha\mathcal H(\pi(\cdot\mid o_t))\}\right],\quad
a=\tfrac12(z_1+1),\quad b=z_2\min\{.44,\sqrt{1-a^2}\}.
\]
\[
r_t=-\ell_t-40(e_v/.05)^2-10e_I^2-2\mathbf1\{e_v>0\text{ or }e_I>0\}.
\]
\[
e_v=[.95-\hat v_{\min}]_++[\hat v_{\max}-1.05]_+,\qquad e_I=[\hat\iota-1]_+.
\]
\referencefont
\begin{tabular}{@{}>{\raggedright\arraybackslash}p{.24\linewidth}>{\raggedright\arraybackslash}p{.69\linewidth}@{}}
$\pi$; $z\in[-1,1]^2$; $a,b$ & Policy, normalized action with components $z_1,z_2$, and dimensionless fleet P/Q participation. Evaluation uses tanh of the Gaussian mean before mapping. \\
$o_t\in\mathbb R^7$ & Three forecast demand multipliers, previous $a,b$, allowed current fraction, remaining horizon/3; all dimensionless. \\
$\mathbb E_\pi,\mathcal H,\alpha=.005$ & Expectation under the policy/environment; differential entropy in normalized action coordinates (natural logarithm); fixed entropy coefficient. \\
$t,\gamma=.98,r_t,\ell_t$ & Training step, discount, reward and engineering cost. Training predicts its voltage-MSE term directly with the fitted response surface. Reward/cost are dimensionless. \\
$e_v,\hat v_{\min},\hat v_{\max}$ & Voltage excess and predicted minimum/maximum voltage, in pu. A hat denotes a fitted response; $[x]_+=\max(x,0)$. \\
$e_I,\hat\iota$ & Dimensionless current excess and predicted maximum line-phase current/rating ratio. Each line-phase uses its own rating; training fixes allowed fraction $\rho=1$. \\
$\mathbf1\{\cdot\}$ & Indicator: one when any stated surrogate violation occurs, zero otherwise. \\
\end{tabular}
\par\textbf{Protocol:} seed 23; 12,000 reward-only transitions; two 64-unit hidden layers; no MPC labels. Fresh network checks remain necessary.
\slidetransition{Which restoration variables enforce service and radiality?}
\sources{Haarnoja et al. (2018), \url{https://proceedings.mlr.press/v80/haarnoja18b.html}. Protocol: \texttt{models/continuous\_sac/training\_protocol.json}; optional \texttt{Continuous\_SAC\_Training.ipynb}.}
\note{Optional historical/detail reference, outside the current 180-minute route. Original notebook IDs and experiment scopes are preserved. Optional. The entropy coefficient is fixed, not automatically learned. The actor and critic optimize the declared surrogate environment. Training return, fit-validation error and fresh physical evaluation are different evidence sets. The action map enforces the device circle but cannot enforce voltage or line-current constraints. The nominal current used in the training penalty is intentional; it makes the reserve case a distribution shift. Spoken bridge: Which restoration variables enforce service and radiality?}
\end{frame}

\begin{frame}[t]{Optional reference: restoration MILP variables and topology}
\setlength{\abovedisplayskip}{3pt}\setlength{\belowdisplayskip}{3pt}\setlength{\parskip}{0pt}
\referencefont
\[
\max\sum_{i\ne0}w_iP_i^d s_i,\quad s_i\leq y_i,\quad
-\overline P_{ij}z_{ij}\leq P_{ij}\leq\overline P_{ij}z_{ij},\quad
\sum_jP_{ji}-\sum_jP_{ij}=P_i^d s_i\quad(i\ne0).
\]
\[
\sum_j f_{ji}-\sum_j f_{ij}=y_i\quad(i\ne0),\quad |f_{ij}|\leq5z_{ij},\quad
\sum_{(i,j)\in E}z_{ij}=\sum_{i\in N}y_i-1.
\]
\referencefont
\renewcommand{\arraystretch}{.98}
\begin{tabular}{@{}>{\raggedright\arraybackslash}p{.25\linewidth}>{\raggedright\arraybackslash}p{.68\linewidth}@{}}
$N,E,i,j$ & Fixture's six buses and candidate edges, with source bus $0$; each edge has one arbitrary reference orientation. Signed flows may oppose it. \\
$s_i,y_i,z_{ij}\in\{0,1\}$ & Load picked up, bus energized, and branch energized. $s_i\leq y_i$ permits an energized transit bus with its load disconnected. \\
$P_i^d\geq0,\overline P_{ij}\geq0,P_{ij}\in\mathbb R$ & Demand, branch capacity and signed active flow, all in kW. Source export equals total picked-up demand. \\
$w_i>0$ & Dimensionless service weight: 10 for the two critical loads, 1 otherwise. The objective is weighted kW; finite weights are not lexicographic priority. \\
$f_{ij}\in\mathbb R$ & Dimensionless connectivity commodity; each energized non-source bus consumes one unit. Source exports $\sum_{i\ne0}y_i$; 5 bounds its total. \\
$y_0=1$, $z_{ij}\leq y_i,y_j$ & Source remains energized; an energized edge requires energized endpoints. Unavailable edges have $z_{ij}=0$. \\
\end{tabular}
\par\textbf{Why radial:} connectivity plus the displayed edge count enforces a tree; edge count alone permits disconnected cycles. A second solve fixes weighted service and minimizes energized branch count.
\slidetransition{How is the residual statistic calibrated?}
\sources{Chen et al. (2018), \url{https://doi.org/10.1109/TPWRS.2017.2720122}; Wang et al. (2020), \url{https://arxiv.org/abs/1912.05185}. This is the six-node active-power fixture, not the feeder catalog.}
\note{Optional historical/detail reference, outside the current 180-minute route. Original notebook IDs and experiment scopes are preserved. Optional. The balance sign follows the chosen edge orientation, with outgoing flow subtracted from incoming flow at non-source buses. Binary y, s and z have different meanings. Edge count alone would permit a disconnected cycle; the commodity constraints exclude it. Capacity is active power here, not a three-phase thermal calculation. The secondary objective is not switch operations from a prior topology. Spoken bridge: How is the residual statistic calibrated?}
\end{frame}

\begin{frame}[t]{Optional reference: WLS, residual covariance and calibration}
\referencefont
\[
\hat x=\arg\min_{x\in\mathbb R^n}(z-Hx)^\top R^{-1}(z-Hx),\quad r=z-H\hat x,
\quad\Omega=R-H(H^\top R^{-1}H)^{-1}H^\top.
\]
\[
r_i^N=|r_i|/\sqrt{\Omega_{ii}},\qquad J=r^\top R^{-1}r,\qquad
J>\chi^2_{m-n,\,1-\alpha}\ \Longrightarrow\ \text{residual alarm}.
\]
\referencefont
\begin{tabular}{@{}>{\raggedright\arraybackslash}p{.23\linewidth}>{\raggedright\arraybackslash}p{.70\linewidth}@{}}
$x,\hat x\in\mathbb R^n$ & True and estimated non-source voltage angles (radians), $n=3$ in the four-bus DC fixture. \\
$z\in\mathbb R^m,H\in\mathbb R^{m\times n}$ & $m=8$ measurements and their linear map: four branch flows, three injections (pu), and one angle (rad). Each row of $H$ maps rad to that channel's units. \\
$R\succ0,\Omega\in\mathbb R^{m\times m}$ & Noise and residual covariance matrices, with channel-dependent squared/cross units. $R$ is diagonal for the stated independent noise. \\
$r,r_i,r_i^N$ & Residual vector, its channel $i$ component in native units, and dimensionless normalized magnitude. Defined only for $\Omega_{ii}>0$. \\
$J,\chi^2_{\nu,1-\alpha},\nu=m-n$ & Dimensionless weighted residual statistic, chi-square quantile and degrees of freedom; $\alpha=.01$, $\nu=5$, threshold 15.0863. \\
$\top$, $R^{-1}$ & Matrix transpose and inverse. Implementation uses weighted least squares/linear solves; the inverse notation states the mathematics. \\
\end{tabular}
\par\smallskip\textbf{Calibration assumptions:} correct linear model, full column rank $\operatorname{rank}(H)=n$, and Gaussian noise with the stated covariance. A residual alarm signals inconsistency; it does not identify malicious intent. Recheck rank after removing a channel.
\slidetransition{Why can a coordinated corruption leave the residual unchanged?}
\sources{Schweppe and Wildes (1970), \url{https://doi.org/10.1109/TPAS.1970.292678}; Handschin et al. (1975), \url{https://doi.org/10.1109/T-PAS.1975.31858}.}
\note{Optional historical/detail reference, outside the current 180-minute route. Original notebook IDs and experiment scopes are preserved. Optional. Distinguish m measurements from n unknown states, and explain whitening through R. The displayed alpha is a nominal false-alarm probability under the stated model, not a measured attack-detection rate. The largest normalized residual is a localization heuristic. Zero residual variance gives an undefined normalized statistic, and a rank-deficient retained set does not support the same estimator. Spoken bridge: Why can a coordinated corruption leave the residual unchanged?}
\end{frame}

\begin{frame}[t]{Optional reference: why a coordinated corruption can be invisible}
\referencefont
\[
z'=z+Hc,\qquad \hat x'=\hat x+c,\qquad z'-H\hat x'=z-H\hat x=r.
\]
\referencefont
\begin{tabular}{@{}>{\raggedright\arraybackslash}p{.25\linewidth}>{\raggedright\arraybackslash}p{.68\linewidth}@{}}
\toprule Symbol / domain & Meaning in the fixed linear WLS model \\
\midrule
$z,z'\in\mathbb R^m$ & Original and corrupted measurement vectors, retaining each channel's native units. \\
$H\in\mathbb R^{m\times n}$ & Known, unchanged full-column-rank measurement matrix; here $m=8,n=3$. \\
$c\in\mathbb R^n$ & Attacker's selected angle shift (rad). The added measurement vector $Hc$ has the correct units for every channel. \\
$\hat x,\hat x'\in\mathbb R^n$ & Original and shifted state estimates (rad); $r$ is their identical residual in measurement units. \\
$\mathcal T,H_{\mathcal T}$ & Index set of protected measurement channels and the corresponding rows of $H$. \\
\bottomrule
\end{tabular}
\par\medskip If protected channels cannot change, a corruption of this form must satisfy $H_{\mathcal T}c=0$. When $H_{\mathcal T}$ has full column rank, only $c=0$ is possible within this attack construction.
\par\medskip\textbf{Threat assumptions:} the attacker knows $H$ and can modify every unprotected channel required by $Hc$. This does not establish that an arbitrary attacker can create the same effect. A low residual does not reject consistency under the stated model; it does not establish the truth of the measurements.
\slidetransition{How does protected information enter constrained control?}
\sources{Liu, Ning and Reiter (2009), \url{https://doi.org/10.1145/1653662.1653666}. Exact fixture: \texttt{src/standard\_formulations.py}, \texttt{run\_wls\_demo}.}
\note{Optional historical/detail reference, outside the current 180-minute route. Original notebook IDs and experiment scopes are preserved. Optional. Substitute the shifted estimate directly into the residual and cancel Hc. This exact cancellation is a property of this linear model with an unchanged matrix, not a blanket statement about nonlinear AC state estimation. The protected-set rank argument rules out this particular Hc shift when the stated channels are trustworthy; it does not authenticate channels or address every attack class. Spoken bridge: How does protected information enter constrained control?}
\end{frame}

\begin{frame}[t]{Optional reference: trusted-data estimation and constrained control}
\referencefont
\[
H_{\mathcal T}=I_3,\quad
b=FB^{-1}e_2,\qquad
\min_{0\leq u\leq .5}\tfrac12u^2
\quad\text{s.t.}\quad |F\hat x_{\mathcal T}+bu|\leq\overline f-.01\mathbf1.
\]
\referencefont
\begin{tabular}{@{}>{\raggedright\arraybackslash}p{.25\linewidth}>{\raggedright\arraybackslash}p{.68\linewidth}@{}}
$\mathcal T,H_{\mathcal T},I_3$ & Three assumed trustworthy angle channels, their measurement matrix, and the $3\times3$ identity. Additional independent angle channels make the toy state observable; this identity is not a row subset of the original eight-channel matrix. \\
$\hat x_{\mathcal T}\in\mathbb R^3$ & Estimate from protected angles, in rad. Three measurements minus three states leave zero residual-test degrees of freedom. \\
$u\in[0,.5]$ & Flexible active injection at bus 2, in pu, balanced by the source. The dimensionless quadratic cost selects the smallest feasible injection. \\
$F\in\mathbb R^{4\times3},B\in\mathbb R^{3\times3}$ & DC maps from non-source angles to branch flows and nodal injections; entries convert rad to pu power. \\
$e_2=(0,1,0)^\top$, $b\in\mathbb R^4$ & Unit selector of bus-2 injection and induced branch-flow change per unit injection. \\
$\overline f=(.35,.18,.35,.35)^\top$ & Absolute branch-flow limits in pu for branches 01, 12, 23, 03. $\mathbf1\in\mathbb R^4$ is an all-ones vector. \\
$F\hat x_{\mathcal T}+bu$, $.01$ & Predicted branch flows and declared planning margin, both in pu; absolute value and inequality act per branch. \\
\end{tabular}
\par\smallskip\textbf{Check after solving:} verify finite values, $u$ bounds and predicted constraints; then evaluate the action against the known synthetic true state. Protected-data control removes the toy overload, subject to the assumed channel trust. The margin is not a calibrated uncertainty bound.
\slidetransition{Which assumptions still limit the scientific claim?}
\sources{Estimation/attack context: \url{https://doi.org/10.1109/TPAS.1970.292678}; \url{https://doi.org/10.1145/1653662.1653666}. Constrained DC teaching response: \texttt{src/standard\_formulations.py}.}
\note{Optional historical/detail reference, outside the current 180-minute route. Original notebook IDs and experiment scopes are preserved. Optional. F and B use the fixture's signed branch-incidence convention; using the same convention in both yields the action sensitivity b. The decision is a scalar DC injection, not feeder-wide AC cyber recovery. Synthetic truth is available to the evaluator, not to an actual operator. The trusted set is an explicit premise, and Llama cannot establish it. Spoken bridge: Which assumptions still limit the scientific claim?}
\end{frame}

\begin{frame}[t]{Study boundaries and the next scientific comparisons}
\referencefont
\begin{tabular}{@{}>{\raggedright\arraybackslash}p{.27\linewidth}>{\raggedright\arraybackslash}p{.66\linewidth}@{}}
\toprule Record & What it measures \\
\midrule
Proposal-complete suite & Six settings $\times$ two supervisors; 24 Llama calls, plus one separate four-call Qwen example. One disclosed development revision; five cost ties and one worse Llama outcome. \\
Initial workshop runs & Continuous SAC/MPC proposals, checks and local Llama fallback choices under both model backends. \\
Language confirmation & 12 authored requests; 8 clear and 4 requiring clarification, with independent semantic scoring and the same downstream tools. \\
SAC training & 12,000 nominal-domain interactions; six authored model/case evaluations. One training seed does not quantify algorithm variability. \\
Standard-method traces & One permitted numerical attempt around continuous P/Q, a six-node MILP or four-bus DC WLS recovery. \\
Historical 32-run corpus & Eight finite-catalog feeder fixtures, two model backends, two supervisors and separate AC reference audits. \\
\bottomrule
\end{tabular}
\par\smallskip\referencefont Language confirmation: correct completion $5/8$ for the parser and $7/8$ for selective Llama; wrong READY interpretations $1/12$ and $2/12$. Numerators count observed outcomes; denominators are clear requests (8) or all requests (12). The restoration subset does not replace the parent-study result.
\par\smallskip\referencefont\textbf{Next comparison:} predeclare unseen cases, repeat model calls and training seeds, match tool/retry budgets, and report incorrect recommendations separately from correct abstentions. Human time and workload need a separate human study.
\slidetransition{Choose one matched comparison to investigate after class.}
\sources{Study design: \texttt{appendix/Agentic\_Advantage\_Study.md}. Recent benchmark context: Mansutti et al. (2026 preprint), \url{https://arxiv.org/abs/2609.00384}. Reading map: \texttt{appendix/Recent\_Agentic\_AI\_Literature.md}.}
\note{Optional historical/detail reference, outside the current 180-minute route. Original notebook IDs and experiment scopes are preserved. Optional. Keep these experiments separate whenever stating a benefit. A successful tool call, correct intent, physical feasibility, good objective and suitable human review are different outcomes. Nothing here establishes field reliability, globally optimal AC control, an authenticated dispatch system, or human time savings. These denominators and limitations are essential to interpreting the observed Llama results in each lesson. Spoken bridge: Choose one matched comparison to investigate after class.}
\end{frame}
